\section{Morfismos universalmente abiertos}
\label{section:IV.14}

Las secciones~\textsection\textsection~14 y~15 están dedicadas al estudio de la noción de morfismo \emph{universalmente abierto} \sref{IV.2.4.2}.
Ya hemos visto \sref{IV.2.4.6} que un morfismo \emph{plano} localmente de presentación finita es universalmente abierto, siendo falsa la recíproca.
En~\textsection~14 examinamos primero las propiedades de las dimensiones de las fibras de un morfismo universalmente abierto $f:X\to Y$; cuando $X$ e $Y$ son localmente noetherianos, $f$ se comporta a este respecto \sref{IV.14.2.1} como un morfismo plano (cf.~\sref{IV.6.1.2}), y es en particular \emph{equidimensional} cuando es localmente de tipo finito y dominante y $X$ es irreducible \sref{IV.14.2.2}.
Recíprocamente, un morfismo equidimensional localmente de tipo finito $f:X\to Y$ es universalmente abierto cuando se supone además que $Y$ es \emph{geométricamente unirramificado}, y en particular cuando $Y$ es \emph{normal} (criterio de Chevalley, \sref{IV.14.4.4}).
Mostramos también que los morfismos universalmente abiertos localmente de tipo finito $f:X\to Y$ (cuando $X$ e $Y$ son
\oldpage[IV-3]{200}
localmente noetherianos e $Y$ es irreducible) admiten ``suficientes'' \emph{cuasisecciones}, es decir, en un entorno de un punto cerrado $x$ de una fibra $f^{-1}(y)$, existe un subpreesquema cerrado $X'$ de $X$ que contiene a $x$ y tal que la restricción $X'\to Y$ de $f$ es un morfismo \emph{cuasifinito} (por tanto con fibras \emph{discretas}) y dominante \sref{IV.14.5.3}.

En~\textsection~15 estudiamos diversas propiedades de las fibras de los morfismos universalmente abiertos $f:X\to Y$, localmente de tipo finito, en particular cuando $X$ e $Y$ son localmente noetherianos.
Obtenemos así, en particular, un criterio para que un punto $x\in X$ pertenezca a una \emph{única} componente irreducible de $X$, en términos de propiedades de la \emph{fibra} $f^{-1}(f(x))$ de $x$: basta que $Y$ sea geométricamente unirramificado (por ejemplo normal) en el punto $f(x)$ y que $f^{-1}(f(x))$ sea geométricamente puntualmente íntegra en el punto $x$ \sref{IV.15.3.3}; si además $Y$ es localmente íntegro en el punto $f(x)$, entonces $f$ es plano en el punto $x$ y $X$ es localmente íntegro en el punto $x$.
Estudiamos también la variación del \emph{número de componentes conexas geométricas} de una fibra $f^{-1}(y)$; por ejemplo, si $f$ es universalmente abierto y \emph{propio}, y las fibras de $f$ son geométricamente reducidas, este número es \emph{localmente constante} \sref{IV.15.5.7}.
Finalmente, cuando $f$ admite una \emph{sección} $g$ (lo que ocurrirá cuando $X$ sea un \emph{esquema en grupos sobre $Y$}) y, para todo $y\in Y$, se denota por $X_y^0$ la componente conexa de la fibra $X_y=f^{-1}(y)$ en el punto $g(y)$ (la ``componente neutra'' en el caso de grupos), se estudia la unión $X^0$ de los $X_y^0$ para $y\in Y$, y se demuestra \sref{IV.15.6.4} que si $f$ es universalmente abierto y las fibras $X_y$ son geométricamente reducidas, entonces $X^0$ es un conjunto \emph{abierto} de $X$.

\subsection{Morfismos abiertos}
\label{subsection:IV.14.1}

\begin{env}[14.1.1]
\label{IV.14.1.1}
Recordemos \sref{IV.1.10.2} que se dice que una aplicación continua $\psi:X\to Y$ es \emph{abierta en un punto} $x\in X$ si la imagen por $\psi$ de \emph{todo} entorno de $x$ en $X$ es un entorno de $\psi(x)$ en $Y$.

Observemos que esto \emph{no} implica que exista un sistema fundamental de entornos de $x$ cuyas imágenes sean \emph{abiertas} en $Y$.
\end{env}

\begin{proposition}[14.1.2]
\label{IV.14.1.2}
Sean $X$, $Y$ dos espacios topológicos, $\psi:X\to Y$ una aplicación continua, $x$ un punto de $X$, $y=\psi(x)$.
\begin{enumerate}[label=(\roman*)]
  \item Si $\psi$ es abierta en $x$, entonces para todo subconjunto $Y'$ de $Y$ que contenga a $\psi(x)$ la restricción $\psi^{-1}(Y')\to Y'$ de $\psi$ es abierta en el punto $x$.
  \item Supongamos que $Y$ es la unión de una familia localmente finita de subconjuntos cerrados $(Y_i)$ y que para todo $i$ tal que $\psi(x)\in Y_i$, la restricción $\psi^{-1}(Y_i)\to Y_i$ de $\psi$ es abierta en el punto $x$; entonces $\psi$ es abierta en el punto $x$.
  \item Sea $\gamma:X'\to X$ una aplicación continua, $x'$ un punto de $X'$; si la aplicación compuesta $\psi\circ\gamma:X'\to Y$ es abierta en el punto $x'$, $\psi$ es abierta en el punto $\gamma(x')$.
\end{enumerate}
\end{proposition}

\begin{proof}
Si $U$ es un entorno de $x$, se tiene $\psi(U\cap\psi^{-1}(Y'))=\psi(U)\cap Y'$; de aquí se sigue inmediatamente (i).
Para demostrar (ii), observemos que existe un entorno $W$ de $\psi(x)$ en $Y$ que solo corta un número finito de los subconjuntos cerrados $Y_i$ que \emph{contienen} a $\psi(x)$; por tanto $W$ es la unión de los $W\cap Y_i$ para estos índices.
Ahora bien, si $U$ es un entorno de $x$ tal que $\psi(U)\cap Y_i$ es un entorno de $\psi(x)$ en $Y_i$, existe un entorno $V\subset W$ de $\psi(x)$ en $Y$ tal que, para todo $i$ tal que $\psi(x)\in Y_i$,
\oldpage[IV-3]{201}
se tiene $V\cap Y_i\subset\psi(U)\cap Y_i$ y, como la unión de los $V\cap Y_i$ para estos índices es $V$, se tiene $V\subset\psi(U)$; por tanto $\psi(U)$ es un entorno de $\psi(x)$ en $Y$.
La afirmación (iii) es trivial.
\end{proof}

\begin{remarks}[14.1.3]
\label{IV.14.1.3}
\begin{enumerate}[label=(\roman*)]
  \item El conjunto $Z$ de los puntos $x\in X$ en los que un morfismo es abierto \emph{no es necesariamente abierto}.
  Por ejemplo, sea $K$ un cuerpo, $A$ el anillo de polinomios $K[S,T]$, $V$ el plano afín $\Spec(A)$, $X$ el subpreesquema cerrado de $V$ que es la ``unión'' de la recta $X_1$ definida por $T=0$ y de la recta $X_2$ definida por $S=0$, es decir, $\Spec(A/\mathfrak{a})$, donde $\mathfrak{a}=AST$; sea $Y=X_2=\Spec(A/AS)$ y sea $f$ la proyección correspondiente a la inyección canónica $K[T]\to A/\mathfrak{a}$; entonces se tiene $Z=X_2$, que no es abierto en $X$.

  \item Sean $X$, $Y$ dos preesquemas noetherianos irreducibles con puntos genéricos $\xi$, $\eta$ respectivamente, $f:X\to Y$ un morfismo \emph{dominante localmente de tipo finito}; entonces $f$ \emph{es abierto en el punto $\xi$}.
  En efecto, podemos evidentemente reducirnos al caso en que $X$ e $Y$ son reducidos (por tanto íntegros) \sref{IV.1.5.4} y afines; en virtud del teorema de planitud genérica \sref{IV.6.9.1}, existe un abierto no vacío $V$ de $Y$ tal que la restricción $f^{-1}(V)\to V$ de $f$ es un morfismo \emph{plano}, y de \sref{IV.2.4.6} concluimos que esta restricción es un morfismo abierto.
  Sin embargo, existen morfismos dominantes de tipo finito $f:X\to Y$ (donde $X$ e $Y$ son \emph{irreducibles}) que no son abiertos (véase por ejemplo \sref[II]{II.8.1}); el conjunto de puntos en los que tal morfismo es abierto no es necesariamente \emph{cerrado} en $X$.

  \item No sabemos si, cuando $X$ e $Y$ son localmente noetherianos y $f:X\to Y$ es un morfismo localmente de tipo finito, el conjunto de puntos de $X$ en los que $f$ es abierto es \emph{localmente construible}.
\end{enumerate}
\end{remarks}

\begin{proposition}[14.1.4]
\label{IV.14.1.4}
Sean $X$, $Y$ dos espacios topológicos, $\psi:X\to Y$ una aplicación continua.
Para todo $y\in Y$, el conjunto de los $x\in\psi^{-1}(y)$ en los que $\psi$ es abierta es un subconjunto cerrado de $\psi^{-1}(y)$.
\end{proposition}

\begin{proof}
En efecto, supongamos que $\psi$ no es abierta en un punto $x\in\psi^{-1}(y)$; existe entonces un entorno $V$ de $x$ en $X$ tal que $\psi(V)$ no es un entorno de $y$; se sigue que para todo $x'\in V\cap\psi^{-1}(y)$, $\psi$ no es abierta en el punto $x'$.
\end{proof}

\begin{remark}[14.1.5]
\label{IV.14.1.5}
Incluso si $X$ e $Y$ son localmente noetherianos y $f$ es un morfismo \emph{finito}, puede ocurrir que $f$ sea abierto en todo punto de una fibra $f^{-1}(y)$, sin que exista un \emph{entorno} de $f^{-1}(y)$ en el que $f$ sea abierto en todos sus puntos.
Tomemos por ejemplo un cuerpo $K$, el anillo de polinomios $A$ igual a $K[\mathrm{T}_1,\mathrm{T}_2,\mathrm{T}_3]$, el espacio afín $V=\Spec(A)$ de dimensión $3$ sobre $K$, y $X=\Spec(A/\mathfrak{a})$, donde $\mathfrak{a}=\mathfrak{b}\mathfrak{c}$, con $\mathfrak{b}=(\mathrm{T}_3)$ y $\mathfrak{c}=(\mathrm{T}_1)+(\mathrm{T}_2-\mathrm{T}_3)$ en $A$, de modo que $X$ es la unión del plano $X_1=\Spec(A/\mathfrak{b})$ (``plano de ecuación $\mathrm{T}_3=0$'') y de la recta $X_2=\Spec(A/\mathfrak{c})$ (``recta de ecuaciones $\mathrm{T}_1=0$, $\mathrm{T}_2=\mathrm{T}_3$~''), que son sus componentes irreducibles.
Tomemos $Y=X_1$ y sea $f:X\to Y$ la proyección correspondiente a la inyección canónica $K[\mathrm{T}_1,\mathrm{T}_2]\to A/\mathfrak{a}$; si $y$ es el punto común a $X_1$ y $X_2$, $f^{-1}(y)$ se reduce a $y$ y $f$ es abierto en este punto, pero no es abierto en ningún punto de $X_2$ situado en un entorno de $y$ y distinto de $y$.
\end{remark}

\begin{env}[14.1.6]
\label{IV.14.1.6}
En lo que sigue, desempeñará un papel esencial el criterio \sref{IV.1.10.3} que, mediante el ``levantamiento de las generalizaciones'', caracteriza los morfismos \emph{localmente de presentación finita} $f:X\to Y$ que son abiertos en un punto $x$: para toda generalización $y'$ de $y=f(x)$, existe $x'\in X$, una generalización de $x$, tal que $y'=f(x')$.
\end{env}

\oldpage[IV-3]{202}
\subsection{Morfismos abiertos y fórmula de la dimensión}
\label{subsection:IV.14.2}

\begin{theorem}[14.2.1]
\label{IV.14.2.1}
Sean $X$, $Y$ dos preesquemas localmente noetherianos, $f:X\to Y$ un morfismo, $x$ un punto de $X$, $y=f(x)$.
Supongamos que $f$ es abierto en los puntos genéricos de las componentes irreducibles de $f^{-1}(y)$ que contienen a $x$.
Entonces se tiene la relación
\begin{equation}
\tag{14.2.1.1}
\label{IV.14.2.1.1}
  \dim(\sh{O}_x)=\dim(\sh{O}_y)+\dim(\sh{O}_x\otimes_{\sh{O}_y}k(y)).
\end{equation}
\end{theorem}

\begin{proof}
Sea $y'$ una generalización cualquiera de $y$ distinta de $y$, y consideremos el subpreesquema cerrado reducido $Y'$ de $Y$ cuyo espacio subyacente es $\overline{\{y'\}}$; ninguna componente irreducible $X'$ de $f^{-1}(Y')$ que contenga a $x$ puede estar contenida en $f^{-1}(y)$: en efecto, si $x'$ es el punto genérico de $X'$, se tendría $f(x')=y\neq y'$, y como $x'$ es su única generalización en $f^{-1}(Y')$, esto contradiría la hipótesis de que $f$ es abierto en $x'$, en virtud del hecho de que $a)$ implique $c)$ en \sref{IV.1.10.3}.
Podemos aplicar, por tanto, \sref{IV.6.1.2} al homomorfismo local de anillos locales noetherianos $\sh{O}_y\to\sh{O}_x$, de donde resulta la conclusión.
\end{proof}

\begin{corollary}[14.2.2]
\label{IV.14.2.2}
Sea $Y$ un preesquema localmente noetheriano, $f:X\to Y$ un morfismo localmente de tipo finito, $x$ un punto de $X$, $y=f(x)$.
Supongamos que $f$ es abierto en los puntos genéricos de las componentes irreducibles de $f^{-1}(y)$ que contienen a $x$.
Entonces $f$ es equidimensional en el punto $x$ en cada uno de los dos casos siguientes:
\begin{enumerate}[label=(\roman*)]
  \item $X$ es irreducible y $f$ es dominante.
  \item Los anillos $\sh{O}_x$ y $\sh{O}_y$ son equidimensionales.
\end{enumerate}
\end{corollary}

\begin{proof}
Para (i), esto se sigue de \sref{IV.14.2.1} y de \sref{IV.13.2.3}.
Para (ii), esto se sigue de \sref{IV.14.2.1} y de \sref{IV.13.3.6}.
\end{proof}

\begin{proposition}[14.2.3]
\label{IV.14.2.3}
Sea $Y$ un preesquema irreducible localmente noetheriano, $\eta$ su punto genérico, $f:X\to Y$ un morfismo localmente de tipo finito, $y$ un punto de $f(X)$, $Z$ una componente irreducible de $f^{-1}(y)$ tal que $f$ es abierto en el punto genérico de $Z$.
Entonces $Z$ está contenida en una componente irreducible $X'$ de $X$ dominante sobre $Y$ y tal que $\dim(Z)=\dim(X'\cap f^{-1}(y))=\dim(X'\cap f^{-1}(\eta))$.
\end{proposition}

\begin{proof}
Esto se sigue, en efecto, de \sref{IV.14.2.1} aplicado en el punto genérico de $Z$, y de \sref{IV.13.2.7}.
\end{proof}

\begin{corollary}[14.2.4]
\label{IV.14.2.4}
Con las notaciones de \sref{IV.14.2.3}, supongamos que $f$ es abierto en todos los puntos de $f^{-1}(y)$.
Para todo $z\in Y$, sea $\mathrm{E}(z)$ el conjunto de las dimensiones de las componentes irreducibles de $f^{-1}(z)$ y $d(z)=\sup\mathrm{E}(z)=\dim(f^{-1}(z))$.
Entonces $\mathrm{E}(y)\subset\mathrm{E}(\eta)$, de donde $d(y)\leq d(\eta)$.
\end{corollary}

\begin{proof}
Esto se sigue inmediatamente de \sref{IV.14.2.3}.
\end{proof}

\begin{corollary}[14.2.5]
\label{IV.14.2.5}
Sea $Y$ un preesquema localmente noetheriano, $f:X\to Y$ un morfismo propio, $y\in Y$ un punto tal que $f$ es abierto en todos los puntos de $f^{-1}(y)$.
Entonces la función $z\mapsto\dim(f^{-1}(z))$ es constante en un entorno de $y$.
\end{corollary}

\begin{proof}
Sean $Y_i$ ($1\leq i\leq n$) los subpreesquemas cerrados reducidos de $Y$ cuyos espacios subyacentes son las componentes irreducibles de $Y$ que contienen a $y$; si $f_i:f^{-1}(Y_i)\to Y_i$ es la restricción de $f$, sabemos que cada uno de los $f_i$ es propio \sref[II]{II.5.4.5} y abierto en todos los puntos de $f^{-1}(y)$ \sref{IV.14.1.2}.
Como la unión de los $Y_i$ es un entorno de $y$ en $Y$,
\oldpage[IV-3]{203}
vemos que podemos reducirnos a demostrar el corolario cuando $Y$ es irreducible; sea $\eta$ su punto genérico.
Se sigue de \sref{IV.14.2.4} que $d(y)\leq d(\eta)$; por otra parte, como $f$ es propio, deducimos de \sref{IV.13.1.5} que $z\mapsto\dim(f^{-1}(z))$ es semicontinua superiormente; en particular $d(y)\geq d(\eta)$, luego $d(y)=d(\eta)$.
Además, existe un entorno $V$ de $y$ tal que $d(z)\leq d(y)$ para todo $z\in V$; pero como todo $z$ es una especialización de $\eta$, se tiene por otra parte $d(z)\geq d(\eta)$, de donde finalmente $d(z)=d(y)$ para $z\in V$.
\end{proof}

\begin{corollary}[14.2.6]
\label{IV.14.2.6}
Sea $Y$ un preesquema localmente noetheriano, $f:X\to Y$ un morfismo abierto de tipo finito, $y$ un punto de $Y$.
Sean $y_i$ ($1\leq i\leq n$) los puntos genéricos de las componentes irreducibles de $Y$ que contienen a $y$.
Entonces, con las notaciones de \sref{IV.14.2.4}:
\begin{enumerate}[label=(\roman*)]
  \item Si para $1\leq i\leq n$ se tiene $\mathrm{E}(y)=\mathrm{E}(y_i)$ (resp.\ $d(y)=d(y_i)$), la función $\mathrm{E}$ (resp.\ $d$) es constante en un entorno de $y$.
  \item Existe un entorno abierto $U$ de $f^{-1}(y)$ tal que la función $z\mapsto\dim(U\cap f^{-1}(z))$ es constante en un entorno de $y$.
\end{enumerate}
\end{corollary}

\begin{proof}
\begin{enumerate}[label=(\roman*)]
  \item El mismo razonamiento que en \sref{IV.14.2.5} muestra que podemos reducirnos a considerar el caso en que $Y$ es irreducible, con punto genérico $\eta$.
  Si $y'$ es una generalización cualquiera de $y$, se tiene entonces (aplicando \sref{IV.14.2.4} a la restricción $f^{-1}(Y')\to Y'$ de $f$, donde $Y'=\overline{\{y'\}}$, y usando \sref{IV.14.1.2}) $\mathrm{E}(y)\subset\mathrm{E}(y')\subset\mathrm{E}(\eta)$ y $d(y)\leq d(y')\leq d(\eta)$; esto muestra que la relación $\mathrm{E}(y)=\mathrm{E}(\eta)$ (resp.\ $d(y)=d(\eta)$) implica $\mathrm{E}(y)=\mathrm{E}(y')$ (resp.\ $d(y)=d(y')$) para toda generalización $y'$ de $y$.
  Ahora bien, en virtud de \sref{IV.9.5.5}, las funciones $\mathrm{E}$ y $d$ son localmente construibles, luego se sigue de $(\mathbf{0}_{\mathrm{III}},~9.2.5)$ aplicado al conjunto de los $z$ tales que $\mathrm{E}(z)=\mathrm{E}(y)$ (resp.\ $d(z)=d(y)$) que este conjunto es un entorno de $y$.

  \item Para cada uno de los $y_i$, $f^{-1}(y_i)$ es un espacio noetheriano puesto que $f$ es de tipo finito; sean $x_{ij}$ ($1\leq j\leq m_i$) los puntos genéricos de sus componentes irreducibles, y pongamos $X_{ij}'=\overline{\{x_{ij}\}}$; mostremos que el complementario $U$ en $X$ de la unión de los $X_{ij}'$ que no cortan a $f^{-1}(y)$ (que evidentemente es un entorno abierto de $f^{-1}(y)$) satisface lo requerido.
  En efecto, la restricción $f|U:U\to Y$ de $f$ es un morfismo abierto de tipo finito \sref[I]{I.6.3.5}; para todo par $(i,j)$ tal que $x_{ij}\in U$, $x_{ij}$ es un punto maximal de $U\cap f^{-1}(y_i)$, y se sigue de \sref{IV.13.1.1}, aplicado a la restricción $X_{ij}'\to Y_i=\overline{\{y_i\}}$ de $f$ (teniendo en cuenta \sref[I]{I.5.2.2} y \sref{IV.1.5.4}), que todas las componentes irreducibles de $X_{ij}'\cap f^{-1}(y)$ tienen dimensión $\geq\dim(X_{ij}'\cap f^{-1}(y_i))$.
  Teniendo en cuenta \sref{IV.14.2.3} aplicado a la restricción $f^{-1}(Y_i)\to Y_i$ de $f$, que es un morfismo abierto \sref{IV.14.1.2}, $f^{-1}(y)$ es, para cada $i$, la unión de las componentes irreducibles de $X_{ij}'\cap f^{-1}(y)$ ($1\leq j\leq m_i$), luego $\dim(f^{-1}(y))\geq\dim(U\cap f^{-1}(y_i))$; pero como $f|U$ es abierto, también se tiene $\dim(f^{-1}(y))\leq\dim(f^{-1}(y_i)\cap U)$ en virtud de \sref{IV.14.2.4}; vemos, por tanto, que se puede aplicar (i) a $f|U$, de donde resulta la conclusión.
\end{enumerate}
\end{proof}

\begin{remark}[14.2.7]
\label{IV.14.2.7}
El ejemplo \sref{IV.13.2.12},~(iii)) muestra que, bajo las hipótesis de \sref{IV.14.2.5} o \sref{IV.14.2.6},~(ii)), no se puede reemplazar en la conclusión la función $z\mapsto\dim(f^{-1}(z))$ por la función $z\mapsto\mathrm{E}(z)$; en efecto, en este ejemplo, $f$ es propio y plano (por tanto universalmente abierto \sref{IV.2.4.6}) y todo entorno del único punto de $f^{-1}(y_0)$ contiene los puntos genéricos de las componentes irreducibles de $f^{-1}(\eta)$.
\end{remark}

\oldpage[IV-3]{204}

\subsection{Morfismos universalmente abiertos}
\label{subsection:IV.14.3}
\label{IV.14.3}

\begin{env}[14.3.1]
\label{IV.14.3.1}
Recordemos \sref{IV.2.4.2} que decir que un morfismo $f:X\to Y$ es universalmente abierto significa que para todo morfismo $g:Y'\to Y$, $f_{(Y')}:X_{(Y')}\to Y'$ es abierto; basta además que así sea cuando $Y'=Y[\mathrm{T}_1,\ldots,\mathrm{T}_e]$, para todo entero $e$ \sref{IV.8.10.2} (y si $Y$ es localmente noetheriano, basta por tanto que así sea para todo $Y'$ localmente noetheriano).
\end{env}

\begin{proposition}[14.3.2]
\label{IV.14.3.2}
Sea $f:X\to Y$ un morfismo de preesquemas.
Si $f$ es universalmente abierto, entonces, para todo morfismo $g:Y'\to Y$, donde $Y'$ es irreducible, la imagen por $f'=f_{(Y')}:X_{(Y')}\to Y'$ de toda componente irreducible de $X_{(Y')}$ es densa en $Y'$.
Recíprocamente, si esta condición se satisface para todo $Y'$ irreducible y todo morfismo de tipo finito $g:Y'\to Y$, y si además $f$ es localmente de presentación finita, entonces $f$ es universalmente abierto.
\end{proposition}

\begin{proof}
En efecto, se sigue de \sref{IV.1.10.4} que si $f$ es universalmente abierto, satisface la condición del enunciado.
Recíprocamente, supongamos que $f$ es localmente de presentación finita, y mostremos que para todo entero $e$, si se pone $Y''=Y[\mathrm{T}_1,\ldots,\mathrm{T}_e]$, $f_{(Y'')}$ es un morfismo abierto.
En efecto, sea $Y'$ un subpreesquema cerrado de $Y''$ cuyo espacio subyacente es un subconjunto cerrado irreducible de $Y''$; el morfismo compuesto $Y'\to Y''\to Y$ es de tipo finito, luego toda componente irreducible de $X_{(Y')}$ domina a $Y'$ por hipótesis; deducimos, por tanto, de \sref{IV.1.10.4} que $f_{(Y'')}$ es un morfismo abierto.
\end{proof}

Esta proposición muestra que la definición \sref[III]{III.4.3.9} coincide, en el caso considerado, con la definición general de los morfismos universalmente abiertos dada en \sref{IV.2.4.2}.

La noción de morfismo abierto en un punto \sref{IV.14.1.1} da lugar del mismo modo a lo siguiente:

\begin{definition}[14.3.3]
\label{IV.14.3.3}
Sea $f:X\to Y$ un morfismo de preesquemas, $x$ un punto de $X$.
Se dice que $f$ es \emph{universalmente abierto en el punto $x$} si, para todo morfismo $g:Y'\to Y$, poniendo $X'=X\times_Y Y'$, el morfismo $f'=f_{(Y')}:X'\to Y'$ es abierto en todo punto $x'$ de $X'$ cuya proyección en $X$ es $x$.
\end{definition}

\begin{remarks}[14.3.3.1]
\label{IV.14.3.3.1}
\begin{enumerate}[label=(\roman*)]
  \item El razonamiento de \sref{IV.8.10.1} muestra (con las mismas notaciones) que si $x$ es un punto de $X$ y $x_\lambda$ es su proyección en $X_\lambda$, entonces, si $f_\mu$ es abierto en el punto $x_\mu$ para todo $\mu\geq\lambda$, $f$ es abierto en el punto $x$; basta limitarse a los abiertos $U$ de $X$ que contienen a $x$ y observar que la hipótesis implica que $f_\mu(U_\mu)$ es un entorno de $f_\mu(x_\mu)$, luego $f(v_\mu^{-1}(U_\mu))$ es un entorno de $f(x)$.
  Deducimos que el enunciado \sref{IV.8.10.2} sigue siendo válido cuando se reemplaza ``universalmente abierto'' por ``universalmente abierto en el punto $x$'', y ``morfismo abierto'' por ``morfismo abierto en todo punto $x_n$ de $X_n$ cuya proyección en $X$ es $x$'': basta limitarse en la demostración a los abiertos $V$ que contienen a $x_n$.

  \item El resultado de \sref{IV.14.1.4} sigue siendo válido para un morfismo $f:X\to Y$, reemplazando ``abierto'' por ``universalmente abierto''.
  En efecto, supongamos que $f$ no es universalmente abierto en un punto $x\in f^{-1}(y)$; existe entonces un morfismo $Y'\to Y$ y un punto $x'\in X'$ que se proyecta sobre $x$, tal que $f'$ no es abierto en el punto $x'$.
\end{enumerate}
\end{remarks}

\oldpage[IV-3]{205}

Ahora bien, si $y'=f'(x')$, la proyección $f'^{-1}(y')\to f^{-1}(y)$ es un morfismo abierto \sref{IV.2.4.10}, y existe, en virtud de \sref{IV.14.1.4}, un entorno $V'$ de $x'$ en $f'^{-1}(y')$ en el que $f'$ no es abierto; por tanto, $f$ no es universalmente abierto en ningún punto de la imagen de $V'$ en $f^{-1}(y)$, que es un entorno de $x$ en $f^{-1}(y)$.

\begin{proposition}[14.3.4]
\label{IV.14.3.4}
\begin{enumerate}[label=(\roman*)]
  \item Sean $f:X\to Y$, $g:Y\to Z$ dos morfismos, $x$ un punto de $X$, $y=f(x)$.
  Si $f$ es universalmente abierto en el punto $x$ y $g$ es universalmente abierto en el punto $y$, entonces $g\circ f$ es universalmente abierto en el punto $x$.
  Recíprocamente, si $g\circ f$ es universalmente abierto en el punto $x$, $g$ es universalmente abierto en el punto $y$.

  \item Si $f:X\to Y$ es un $S$-morfismo universalmente abierto en el punto $x\in X$, entonces, para todo cambio de base $S'\to S$, $f_{(S')}:X_{(S')}\to Y_{(S')}$ es universalmente abierto en todo punto de $X_{(S')}$ situado sobre $x$.

  \item Para que $f:X\to Y$ sea universalmente abierto en el punto $x\in X$, es necesario y suficiente que $f_{\mathrm{red}}$ lo sea.

  \item Sea $f:X\to Y$ un morfismo localmente de presentación finita, $x$ un punto de $X$, $y=f(x)$; pongamos $Y_1=\Spec(\sh{O}_y)$, $X_1=X\times_Y Y_1$, $f_1=f_{(Y_1)}$.
  Para que $f$ sea universalmente abierto en el punto $x$, es necesario y suficiente que $f_1$ lo sea (recordemos \sref[I]{I.3.6.5} que $X_1$ se identifica canónicamente con un subespacio de $X$).
\end{enumerate}
\end{proposition}

\begin{proof}
En efecto, (ii) es una consecuencia evidente de la definición \sref{IV.14.3.3}; de la definición se sigue también que para demostrar la afirmación (i), basta hacerlo suprimiendo en todas partes la palabra ``universalmente'', y esto se sigue entonces de \sref{IV.14.1.2}.
La afirmación (iii) se sigue de (ii) y del hecho de que el morfismo canónico $X_{\mathrm{red}}\to X$ es sobreyectivo.
Finalmente, la condición (iv) es trivialmente necesaria.
Por otra parte, si se satisface y si $g:Y'\to Y$ es un morfismo cualquiera, pongamos $X'=X_{(Y')}$, $f'=f_{(Y')}$, y sea $x'$ un punto de $X'$ situado sobre $x$; entonces $f'$ es localmente de presentación finita, y para ver que es abierto en el punto $x'$, basta aplicar el criterio {\EGAReadLink{IV.1.10.3}{(1.10.3, $b)$)}}.
Sea $y'=f'(x')$, $Y_1'=\Spec(\sh{O}_{y'})$, $X_1'=X'\times_{Y'}Y_1'$; como $g(y')=y$, el morfismo compuesto $Y_1'\to Y'\xrightarrow{g}Y$ se factoriza como $Y_1'\to Y_1\to Y$ \sref[I]{I.2.4.4}, luego $X_1'=X_1\times_{Y_1}Y_1'$ y $f_1'=(f_1)_{(Y_1')}$; la conclusión se sigue entonces inmediatamente de la hipótesis de que $f_1'$ es abierto en el punto $x'$ y de \sref{IV.1.10.3}.
\end{proof}

\begin{proposition}[14.3.5]
\label{IV.14.3.5}
Sean $X$, $Y$ dos preesquemas, $f:X\to Y$ un morfismo, $x$ un punto de $X$.
Sea $(Y_i)_{1\leq i\leq n}$ una familia localmente finita de subpreesquemas cerrados de $Y$ tal que el espacio $Y$ es la unión de los $Y_i$, y supongamos que para todo $i$ tal que $f(x)\in Y_i$, la restricción $f^{-1}(Y_i)\to Y_i$ de $f$ es universalmente abierta en el punto $x$; entonces $f$ es universalmente abierto en el punto $x$.
\end{proposition}

\begin{proof}
Teniendo en cuenta la definición, esto se sigue de la proposición análoga {\hyperlink{ega4.c04.prop.14-1-2}{(14.1.2, (ii))}} para los morfismos abiertos en un punto.
\end{proof}

\begin{proposition}[14.3.6]
\label{IV.14.3.6}
Sea $Y$ un preesquema localmente noetheriano, $f:X\to Y$ un morfismo localmente de tipo finito.
Para que $f$ sea universalmente abierto en un punto $x\in X$, es necesario y suficiente que se satisfaga la condición siguiente: para todo morfismo $g:Y'\to Y$, donde $Y'=\Spec(A)$ es el espectro de un anillo de valoración discreta, tal que la imagen $g(y')$ del punto cerrado $y'$ de $Y'$ sea igual a $y=f(x)$, y para todo punto $x'\in X'=X\times_Y Y'$ cuyas proyecciones sobre $X$ e $Y'$ sean respectivamente $x$ e $y'$, existe una generalización $z'$ de $x'$ en $X'$ cuya proyección en $Y'$ es el punto genérico de $Y'$ (en otras palabras, \emph{existe una componente irreducible de $X'$ que contiene a $x'$ y domina a $Y'$}).

\oldpage[IV-3]{206}

Además, en la condición precedente, podemos limitarnos al caso en que $A$ es completo, su cuerpo residual es algebraicamente cerrado y $x'$ es racional sobre $k(y')$.
\end{proposition}

\begin{proof}
Si $Y'$ es como en el enunciado, la necesidad de la condición se sigue del hecho de que $f'$ debe ser abierto en el punto $x'$, y del criterio \sref{IV.1.10.3}.
Para ver que la condición es suficiente, consideremos un morfismo de tipo finito $g:Y''\to Y$\footnote{La fuente impresa tiene aquí $g:Y''\to U$, aunque $U$ no está definido; el producto fibrado posterior es sobre $Y$.}, y sean $X''=X\times_Y Y''$, $f''=f_{(Y'')}:X''\to Y''$, y $x''$ un punto de $X''$ situado sobre $x$.
Pongamos $y''=f''(x'')$, y sea $t$ una generalización de $y''$ en $Y''$, distinta de $y''$.
Como $Y''$ es localmente noetheriano, se sigue de \sref[II]{II.7.1.9} y $(\mathbf{0}_{\mathrm{III}},~10.3.1)$ que existe un esquema $Y'=\Spec(A)$, donde $A$ es un anillo de valoración discreta completo cuyo cuerpo residual es una clausura algebraica de $k(x'')$, y un morfismo $g:Y'\to Y''$ tal que, si $s$ e $y'$ son respectivamente el punto genérico y el punto cerrado, se tiene $g(s)=t$ y $g(y')=y''$.
Existe entonces un punto $x'$ de $X'=X\times_Y Y'=X''\times_{Y''}Y'$ cuyas proyecciones en $X''$ e $Y'$ son respectivamente $x''$ e $y'$, y que es racional sobre $k(y')$ \sref[I]{I.3.4.9}.
La hipótesis implica que existe una generalización $z'$ de $x'$ en $X'$ cuya proyección en $Y'$ es $s$; si $z''$ es la proyección de $z'$ en $X''$, $z''$ es una generalización de $x''$ y su proyección en $Y''$ es $t$; concluimos, por tanto, de \sref{IV.1.10.3} que $f''$ es abierto en el punto $x''$, luego $f$ es universalmente abierto en el punto $x$ {\hyperlink{ega4.c04.subsec.14-3-3-1}{(14.3.3.1, (i))}}.
\end{proof}

\begin{corollary}[14.3.7]
\label{IV.14.3.7}
Con las notaciones de \sref{IV.14.3.6}:
\begin{enumerate}[label=(\roman*)]
  \item Dado un punto $y\in Y$, para que $f$ sea universalmente abierto en todos los puntos de $f^{-1}(y)$, es necesario y suficiente que para todo morfismo $g:Y'\to Y$, donde $Y'$ es el espectro de un anillo de valoración discreta y donde la imagen $g(y')$ del punto cerrado $y'$ de $Y'$ es $y$, toda componente irreducible de $X'=X\times_Y Y'$ domine a $Y'$.
  \item Para que $f$ sea universalmente abierto, es necesario y suficiente que para todo morfismo $g:Y'\to Y$, donde $Y'$ es el espectro de un anillo de valoración discreta, toda componente irreducible de $X'=X\times_Y Y'$ domine a $Y'$.
\end{enumerate}
\end{corollary}

\begin{proof}
Está claro que basta demostrar (i); la necesidad de (i) se sigue de \sref{IV.14.3.2} y su suficiencia de \sref{IV.14.3.6}.
\end{proof}

\begin{proposition}[14.3.8]
\label{IV.14.3.8}
Sea $Y$ un preesquema localmente noetheriano, irreducible, regular y de dimensión $1$ (por ejemplo, el espectro de un anillo de Dedekind), $f:X\to Y$ un morfismo localmente de tipo finito, $y$ un punto de $Y$.
Las condiciones siguientes son equivalentes:
\begin{enumerate}[label=\emph{\alph*)}]
  \item $f_{\mathrm{red}}$ es plano en todo punto de $f^{-1}(y)$.
  \item $f$ es universalmente abierto en un entorno de $f^{-1}(y)$.
  \item $f$ es abierto en un entorno de $f^{-1}(y)$.
  \item Toda componente irreducible de $X$ que corta a $f^{-1}(y)$ domina a $Y$.
\end{enumerate}
\end{proposition}

\begin{proof}
Como $f_{\mathrm{red}}$ es localmente de tipo finito \sref{IV.1.3.4}, $a)$ implica que $f_{\mathrm{red}}$ es plano en un entorno de $f^{-1}(y)$ \sref{IV.11.1.1}, y basta aplicar \sref{IV.2.4.6} en tal entorno para ver que $a)$ implica $b)$.
La implicación $b)\Rightarrow c)$ es trivial, y la implicación $c)\Rightarrow d)$ se sigue de \sref{IV.1.10.4} aplicado a un entorno de $f^{-1}(y)$.
Queda por ver que $d)$ implica $a)$.
Podemos evidentemente, en virtud de \sref{IV.1.3.4}, reducirnos al caso en que $X$ es reducido.
Como además la cuestión es local en $X$ y en $Y$, podemos suponer que $Y=\Spec(A)$ y $X=\Spec(B)$
\oldpage[IV-3]{207}
afines; si $X_i$ ($1\leq i\leq n$) son los subpreesquemas cerrados íntegros de $X$ cuyos espacios subyacentes son las componentes irreducibles de $X$, entonces, para todo $x\in X$, $\sh{O}_{X,x}$ es un subanillo del producto directo de los $\sh{O}_{X_i,x}$; si $y=f(x)$, basta mostrar que cada $\sh{O}_{X_i,x}$ es un $\sh{O}_y$-módulo \emph{sin torsión}, pues entonces lo mismo será cierto para $\sh{O}_{X,x}$; como por hipótesis $\sh{O}_y$ es un anillo local regular de dimensión $1$, es decir, un anillo de valoración discreta \sref[II]{II.7.1.6}, se sigue entonces de $(\mathbf{0}_1,~6.3.4)$ que $\sh{O}_{X,x}$ es un $\sh{O}_y$-módulo \emph{plano}.
Pero si $X_i=\Spec(B_i)$, donde $B_i$ es un anillo íntegro, la hipótesis implica que el homomorfismo $A\to B_i$ es \emph{inyectivo} \sref[I]{I.1.2.7}; luego $B_i$ es un $A$-módulo sin torsión, y \emph{a fortiori} $\sh{O}_{X_i,x}$ es un $\sh{O}_y$-módulo sin torsión.
\end{proof}

\begin{remarks}[14.3.9]
\label{IV.14.3.9}
\begin{enumerate}[label=(\roman*)]
  \item En el enunciado de \sref{IV.14.3.8}, no se puede prescindir de la hipótesis de que $Y$ sea \emph{regular}.
  Con las notaciones de \sref{IV.11.7.5}, tomemos en efecto $Y=\Spec(A)$, $X=\Spec(\overline{A})$, de modo que $f:X\to Y$ es un morfismo \emph{finito sobreyectivo}; como $A$ es un anillo local íntegro de dimensión $1$, al igual que $\overline{A}$, se sigue inmediatamente de \sref{IV.1.10.4} que el morfismo $f$ es \emph{abierto}.
  Sin embargo, $f$ no es universalmente abierto (ni \emph{a fortiori} plano) \sref{IV.11.7.5}.
  Se podría dar un ejemplo análogo tomando como $Y$ el esquema local en un ``punto doble ordinario'' de una curva algebraica, y como $X$ la normalización de $Y$.

  \item El ejemplo de {\hyperlink{ega4.c04.item.14-1-3}{(14.1.3, (i))}} muestra que el conjunto de puntos de $X$ en los que un morfismo es universalmente abierto no es necesariamente abierto, puesto que el morfismo $f$ de ese ejemplo es universalmente abierto en todo punto en el que es abierto.
  El ejemplo $f:X\to Y$ considerado anteriormente en (i) muestra asimismo que el conjunto de puntos en los que un morfismo es universalmente abierto no es necesariamente cerrado, pues es inmediato que $f$ es universalmente abierto en todos los puntos de $X$ salvo uno (de hecho, en ellos es incluso un isomorfismo local).
  Sería interesante saber si el conjunto de puntos en los que un morfismo es universalmente abierto es localmente construible.
\end{enumerate}
\end{remarks}

Las dos proposiciones siguientes nos fueron señaladas por M.\ Artin:

\begin{lemma}[14.3.10]
\label{IV.14.3.10}
Sea $A$ un anillo de valoración (no necesariamente discreta), $k$ su cuerpo residual, $K$ su cuerpo de fracciones, y pongamos $S=\Spec(A)$.
Sea $X$ un preesquema irreducible, $f:X\to S$ un morfismo dominante de tipo finito; sea $X_0=X\otimes_A k$ (resp.\ $X_1=X\otimes_A K$) la fibra de $f$ en el punto cerrado (resp.\ en el punto genérico) de $S$.
Entonces, si $X_0\neq\varnothing$, se tiene $\dim(X_0)=\dim(X_1)$.
\end{lemma}

\begin{proof}
Podemos reducirnos al caso en que $X$ es afín, reemplazando si es necesario $X$ por un abierto afín que contenga un punto genérico de una componente irreducible de $X_0$ de dimensión máxima.
Sea $n=\dim(X_0)\geq 0$; usando \sref{IV.4.1.1.3}, se sigue de \sref{IV.13.3.1.1} que existe un entorno $U$ de $X_0$ en $X$ y un $S$-morfismo cuasifinito $g:U\to S[\mathrm{T}_1,\ldots,\mathrm{T}_n]=Z$ tal que la restricción $g_0:U_0=U\cap X_0\to Z_0=\Spec(k[\mathrm{T}_1,\ldots,\mathrm{T}_n])$ es finita y \emph{sobreyectiva}.
Por el cambio de base $\Spec(K)\to S$ y la restricción al subconjunto abierto $U_1=U\cap X_1$ de $X_1$, deducimos de $g$ un morfismo cuasifinito $g_1:U_1\to Z_1=\Spec(K[\mathrm{T}_1,\ldots,\mathrm{T}_n])$.
Como $U_1$ es denso en $X_1$ \sref{IV.4.1.1.3}, se tiene $\dim(U_1)=\dim(X_1)$; si podemos demostrar que el morfismo $g_1$ es \emph{dominante}, la proposición quedará establecida, pues entonces $\dim(U_1)\leq\dim(Z_1)=n$ y, por otra parte, $\dim(X_1)\geq n$ por \sref{IV.4.1.2}.
Supongamos lo contrario; existiría entonces un polinomio $\mathrm{F}_1\neq 0$ en $K[\mathrm{T}_1,\ldots,\mathrm{T}_n]$ tal que $g_1(U_1)\subset\mathrm{V}(\mathrm{F}_1)$.
Si $\omega$ es una valoración sobre $K$ asociada a $A$, y si $(a_\alpha)$ es la familia de coeficientes de $\mathrm{F}_1$, después de multiplicar $\mathrm{F}_1$ por un elemento no nulo de $K$, podemos suponer que $\inf_\alpha(\omega(a_\alpha))=0$; en otras palabras, $\mathrm{F}_1$ procede de un polinomio $\mathrm{F}\in A[\mathrm{T}_1,\ldots,\mathrm{T}_n]$ (con el que se identifica) tal que la imagen $\mathrm{F}_0$ de $\mathrm{F}$ en $k[\mathrm{T}_1,\ldots,\mathrm{T}_n]$ es \emph{no nula}.
Consideremos entonces el subconjunto cerrado $\mathrm{V}(\mathrm{F})$ de $Z$; se tiene $\mathrm{V}(\mathrm{F})\cap Z_1=\mathrm{V}(\mathrm{F}_1)$, y como $U_1$
\oldpage[IV-3]{208}
es denso en $U$ (pues contiene el punto genérico de $X$), $g(U)\subset\mathrm{V}(\mathrm{F})$; en particular, $g_0(U_0)\subset\mathrm{V}(\mathrm{F})\cap Z_0=\mathrm{V}(\mathrm{F}_0)$; pero como $\mathrm{F}_0\neq 0$, $\mathrm{V}(\mathrm{F}_0)$ es un subconjunto cerrado de $Z_0$ distinto de $Z_0$, y se llega a una contradicción.
\end{proof}

\begin{proposition}[14.3.11]
\label{IV.14.3.11}
Sea $f:X\to S$ un morfismo de tipo finito, $(g_\lambda)_{\lambda\in L}$ una familia de morfismos universalmente abiertos $g_\lambda:Y_\lambda\to S$, y para todo $\lambda\in L$, sea $u_\lambda:Y_\lambda\to X$ un $S$-morfismo.
Para todo $s\in S$, pongamos $X_s=X\times_S\Spec(k(s))$, $(Y_\lambda)_s=Y_\lambda\times_S\Spec(k(s))$, y sea $(u_\lambda)_s:(Y_\lambda)_s\to X_s$ el morfismo $u_\lambda\times 1$ inducido a partir de $u_\lambda$ por cambio de base.
Sea $Z(s)$ la adherencia en $X_s$ de la unión de los conjuntos $(u_\lambda)_s((Y_\lambda)_s)$, y pongamos $d(s)=\dim(Z(s))$.
Entonces, para toda generalización $s'$ de un punto $s\in S$, se tiene $d(s)\leq d(s')$.
\end{proposition}

\begin{proof}
Sabemos por \sref[II]{II.7.1.4} que existe un anillo de valoración $A'$ y un morfismo $h:S'=\Spec(A')\to S$ tal que, si $a$ (resp.\ $b$) es el punto cerrado (resp.\ el punto genérico) de $S'$, se tiene $h(a)=s$, $h(b)=s'$.
Además, la proyección $p:X_s\otimes_{k(s)}k(a)\to X_s$ es sobreyectiva y abierta \sref{IV.2.4.10}, y hace así de $X_s$ un espacio cociente de $X_s\otimes_{k(s)}k(a)$ por una relación de equivalencia abierta; por consiguiente, para todo subconjunto $M$ de $X_s$, $p^{-1}(\overline{M})=\overline{p^{-1}(M)}$ (Bourbaki, \emph{Topología general}, cap.~I, 4.ª ed., \S5, n.º~3, prop.~7).
Se razona del mismo modo para $X_{s'}$; teniendo en cuenta \sref[I]{I.3.4.8}, \sref{IV.4.2.7}, y el hecho de que los $g_\lambda$ son universalmente abiertos, vemos que basta demostrar la proposición después del cambio de base $S'\to S$.
Supongamos, por tanto, que $S'=S$, que $s$ es el punto cerrado y que $s'$ es el punto genérico de $S$.
La hipótesis de que $g_\lambda$ es abierto implica que toda componente irreducible de $Y_\lambda$ domina a $S$ \sref{IV.1.10.4}, luego su punto genérico es un punto maximal de $(Y_\lambda)_{s'}$; si $Z$ denota la adherencia en $X$ de $Z(s')$, se tiene por tanto $u_\lambda(Y_\lambda)\subset Z$, y por consiguiente $(u_\lambda)_s((Y_\lambda)_s)\subset Z_s=Z\cap X_s$; en otras palabras, $Z(s)\subset Z_s$, de donde $\dim(Z(s))\leq\dim(Z_s)$.
Pero aplicando \sref{IV.14.3.10} a un subpreesquema reducido de $X$ cuyo espacio subyacente sea una componente irreducible de $Z$, se obtiene $\dim(Z_s)\leq\dim(Z_{s'})$ (la desigualdad proviene del hecho de que puede haber componentes irreducibles de $Z$ que no corten a $X_s$).
Por otra parte, como $Z(s')$ es cerrado en $X_{s'}$ por definición, se tiene $Z_{s'}=Z(s')$, lo que termina la demostración.
\end{proof}

\begin{remark}[14.3.12]
\label{IV.14.3.12}
El caso considerado por M.\ Artin era aquel en que $Y_\lambda=S$ para todo $\lambda$, en otras palabras, el caso en que $(u_\lambda)$ es una familia de $S$-\emph{secciones} de $X$.
Otro caso útil es aquel en que la familia $(u_\lambda)$ se reduce a un único elemento; siempre se puede reducir a este caso considerando el preesquema $Y$ que es la suma de los $Y_\lambda$ y los morfismos $g:Y\to S$ y $u:Y\to X$ cuyas restricciones a cada $Y_\lambda$ son respectivamente $g_\lambda$ y $u_\lambda$.
\end{remark}

\begin{proposition}[14.3.13]
\label{IV.14.3.13}
Sea $f:X\to Y$ un morfismo localmente de tipo finito, $y$ un punto de $Y$, $x$ un punto maximal de la fibra $X_y=X\times_Y\Spec(k(y))$.
Consideremos las condiciones siguientes:
\begin{enumerate}[label=\emph{\alph*)}]
  \item $f$ es universalmente abierto en el punto $x$ (o incluso en todo punto de la componente irreducible de $X_y$ de punto genérico $x$ {\hyperlink{ega4.c04.subsec.14-3-3-1}{(14.3.3.1, (ii))}}).
  \item Para toda componente irreducible $Y_0$ de $Y$ que contenga a $y$, existe una componente irreducible $Z$ de $X$ que contiene a $x$, dominante sobre $Y_0$ y tal que $\dim_x(X_y)=\dim_x(Z\cap X_y)=\dim(Z\cap X_\eta)$, donde $\eta$ es el punto genérico de $Y_0$ (lo que implica que $Z$ es \emph{equidimensional} sobre $Y_0$ en el punto $x$ \sref{IV.13.2.2}).
\end{enumerate}

\oldpage[IV-3]{209}

\begin{enumerate}[label=\emph{\alph*$'$)}]
  \item[b$'$)] Para todo entorno abierto $U$ de $x$ en $X$ y toda generalización $y'$ de $y$, se tiene $\dim(U\cap X_{y'})\geq\dim_x(U\cap X_y)$.
\end{enumerate}

Entonces se tienen las implicaciones $a)\Rightarrow b)\Leftrightarrow b')$.
\end{proposition}

\begin{proof}
Para mostrar que $b)$ implica $b'$), basta observar que $y'$ pertenece a una componente irreducible $Y_0$ de $Y$ que contiene a $y$ y cuyo punto genérico es $\eta$; tomando $Z$ como en $b)$ y observando que el punto genérico de $Z$ (que es también el de $Z\cap X_\eta$ $(\mathbf{0}_1,~2.1.8)$) está contenido en $U$, se tiene $\dim(U\cap X_{y'})\geq\dim(U\cap Z\cap X_{y'})$, y, en virtud de \sref{IV.13.1.6}, $\dim(U\cap Z\cap X_{y'})\geq\dim(Z\cap X_\eta)$; de donde resulta la afirmación, ya que $\dim_x(U\cap X_y)=\dim_x(X_y)$ \sref{IV.4.1.1.3}.

Para demostrar que $b')$ implica $b)$, se puede reemplazar primero $X$ por $f^{-1}(Y_0)$, y suponer por tanto $Y_0=Y$; podemos reducirnos al caso en que $X$ e $Y$ son afines, puesto que $\dim_x(U\cap X_y)=\dim_x(X_y)$ \sref{IV.4.1.1.3}.
Las componentes irreducibles $W_i$ del preesquema noetheriano $X_\eta$ son entonces finitas en número, y el complementario de la unión de las $\overline{W}_i$ (adherencias \emph{en} $X$) que no contienen a $x$ es un entorno abierto $V$ de $x$.
Reemplazando $X$ por $V$, podemos suponer por tanto que $x\in\overline{W}_i$ para todo $i$ (la hipótesis $b'$) implica $\dim(V\cap X_\eta)\geq 0$, luego $x$ pertenece a $\overline{W}_i$ para al menos un $i$).
Además, las $\overline{W}_i$ son exactamente las componentes irreducibles de $X$ que dominan a $Y$ $(\mathbf{0}_1,~2.1.8)$; si se tuviera, para cada una de estas componentes $Z$, $\dim(Z\cap X_\eta)<\dim_x(X_y)$, se concluiría $\dim(X_\eta)<\dim_x(X_y)$, en contradicción con la hipótesis $b'$).
La relación $\dim_x(Z\cap X_y)=\dim(Z\cap X_\eta)$ se sigue entonces de \sref{IV.13.1.6}.

Queda por mostrar que $a)$ implica $b'$).
Teniendo en cuenta \sref[II]{II.7.1.4} y la invariancia de las hipótesis y de la conclusión por cambio de base (en virtud de \sref{IV.4.2.7}), podemos reducirnos al caso en que $Y$ es el espectro de un anillo de valoración, con punto cerrado $y$ y punto genérico $y'$, y en que $U=X$.
La hipótesis de que $f$ es abierto en el punto $x$ implica que existe una componente irreducible $Z$ de $X$ que contiene a $x$ y domina a $Y$ \sref{IV.1.10.3}.
Aplicando \sref{IV.14.3.10} a un entorno de $x$ en $Z$, concluimos que $\dim(Z\cap X_{y'})=\dim_x(Z\cap X_y)$; pero como $x$ es maximal en $X_y$, $Z\cap X_y$ contiene la componente irreducible de $X_y$ de punto genérico $x$, luego $\dim_x(Z\cap X_y)=\dim_x(X_y)$; por otra parte, se tiene $\dim(X_{y'})\geq\dim(Z\cap X_{y'})$, lo que termina la demostración de $b'$).
\end{proof}

\begin{remark}[14.3.14]
\label{IV.14.3.14}
No sabemos si, en \sref{IV.14.3.13}, la conclusión sigue siendo válida cuando se reemplaza la hipótesis $a)$ por la hipótesis más débil de que $f$ es \emph{abierto} en el punto $x$.
Sin embargo, se podría mostrar fácilmente que bastaría tratar el caso en que $Y$ es el espectro de un anillo local íntegro, cuyo punto genérico es aislado, y en que $X$ es un subpreesquema cerrado del fibrado vectorial $Y[\mathrm{T}]$.
\end{remark}
\subsection{Criterio de Chevalley para los morfismos universalmente abiertos}
\label{subsection:IV.14.4}

\begin{theorem}[14.4.1]
\label{IV.14.4.1}
Sea $f:X\to Y$ un morfismo localmente de tipo finito, $y$ un punto de $Y$, $x$ un punto maximal de la fibra $X_y=X\times_Y\Spec(k(y))$.
Supongamos que $y$ es geométricamente unirramificado.
Entonces las condiciones siguientes son equivalentes:
\begin{enumerate}[label=\emph{\alph*)}]
  \item $f$ es universalmente abierto en $x$ (o, equivalentemente, en todo punto de la componente irreducible de $X_y$ cuyo punto genérico es $x$ {\hyperlink{ega4.c04.subsec.14-3-3-1}{(14.3.3.1, (iii))}}).
\oldpage[IV-3]{210}
  \item Si $Y_0$ es la única componente irreducible de $Y$ que contiene a $y$, y $\eta$ es su punto genérico, existe una componente irreducible $Z$ de $X$ que contiene a $x$, domina a $Y_0$ y satisface $\dim_x(Z\cap X_y)=\dim(Z\cap X_\eta)$ (lo que significa que $Z$ es equidimensional sobre $Y_0$ en el punto $x$ \sref{IV.13.2.2}).
  \item[\emph{b$'$)}] Para todo entorno abierto $U$ de $x$ en $X$ y toda generalización $y'$ de $y$, se tiene $\dim(U\cap X_{y'})\geq\dim_x(U\cap X_y)$.
\end{enumerate}
Si además $Y$ es localmente noetheriano, estas condiciones son también equivalentes a la siguiente:
\begin{enumerate}[label=\emph{\alph*)}]
  \item[\emph{c)}] $f$ es abierto en $x$.
\end{enumerate}
\end{theorem}

Observemos primero que, como $(\sh{O}_y)_{\mathrm{red}}$ es íntegro, $y$ pertenece a una sola componente irreducible $Y_0$ de $Y$.
El hecho de que \emph{b)} y \emph{b$'$)} sean equivalentes y que \emph{a)} implique \emph{b$'$)} resulta de \sref{IV.14.3.13}; por otra parte, si $Y$ es localmente noetheriano, hemos visto en \sref{IV.14.2.3} que \emph{c)} implica \emph{b)}.
Queda por demostrar que, cuando $y$ es geométricamente unirramificado, \emph{b)} implica \emph{a)}.

\begin{lemma}[14.4.1.1]
\label{IV.14.4.1.1}
Sea $Y$ un preesquema, $Y'=Y[\mathrm{T}_1,\ldots,\mathrm{T}_n]$, $y'$ un punto de $Y'$ y $y$ su imagen en $Y$.
Si $Y$ es geométricamente unirramificado en el punto $y$, entonces $Y'$ es geométricamente unirramificado en el punto $y'$.
\end{lemma}

En efecto, para todo $y\in Y$, $\Spec(\kappa(y)[\mathrm{T}_1,\ldots,\mathrm{T}_n])$ es geométricamente regular sobre $\kappa(y)$ \sref[0]{0.17.3.7}, el morfismo estructural $Y'\to Y$ es liso \sref{IV.6.8.1}, y basta aplicar \sref{IV.11.3.14}.

\begin{lemma}[14.4.1.2]
\label{IV.14.4.1.2}
Sea $A$ un anillo local íntegro y unirramificado, $B$ un anillo íntegro que contiene a $A$ y es entero sobre él, y $\mathfrak{n}$ un ideal primo de $B$ situado sobre el ideal maximal $\mathfrak{m}$ de $A$.
Entonces el morfismo $\Spec(B_\mathfrak{n})\to\Spec(A)$ es sobreyectivo; dicho de otro modo, para todo ideal primo $\mathfrak{p}$ de $A$ existe un ideal primo $\mathfrak{q}$ de $B$ tal que $\mathfrak{q}\subset\mathfrak{n}$ y $\mathfrak{q}\cap A=\mathfrak{p}$.
\end{lemma}

Sea $K$ (resp.\ $L$) el cuerpo de fracciones de $A$ (resp.\ $B$), $A'$ el cierre entero de $A$, y $B'$ el subanillo de $L$ generado por $A'$ y $B$, de modo que tenemos un diagrama conmutativo de inyecciones canónicas
\[
  \xymatrix{
    B \ar[r] & B' \\
    A \ar[u] \ar[r] & A' \ar[u]
  }
\]
Como $B'$ es entero sobre $B$, existe un ideal primo $\mathfrak{n}'$ de $B'$ tal que $\mathfrak{n}'\cap B=\mathfrak{n}$ (Bourbaki, \emph{Alg.\ comm.}, cap.~V, \S~2, n\textsuperscript{o}~1, teo.~1), y (por la misma razón) $\Spec(A')\to\Spec(A)$ es sobreyectivo.
Por otra parte, como $A$ es unirramificado, $A'$ es un anillo \emph{local}, de modo que $\mathfrak{n}'\cap A'$, que está sobre el ideal maximal $\mathfrak{m}$ de $A$, es necesariamente igual al único ideal maximal $\mathfrak{m}'$ de $A'$.
En virtud del segundo teorema de Cohen--Seidenberg (\emph{loc.\ cit.}, \S~2, n\textsuperscript{o}~4, teo.~3), el morfismo $\Spec(B'_{\mathfrak{n}'})\to\Spec(A')$ es sobreyectivo, y lo mismo vale para el compuesto $\Spec(B'_{\mathfrak{n}'})\to\Spec(A')\to\Spec(A)$; pero este morfismo es también el compuesto $\Spec(B'_{\mathfrak{n}'})\to\Spec(B_\mathfrak{n})\to\Spec(A)$, por lo que el morfismo $\Spec(B_\mathfrak{n})\to\Spec(A)$ es sobreyectivo.

Establecidos estos lemas, volvamos a la demostración de la implicación $b)\Rightarrow a)$ en \sref{IV.14.4.1}.
En virtud de {\hyperlink{ega4.c04.subsec.14-3-3-1}{(14.3.3.1, (i))}}, basta demostrar que, para todo entero $n\geq 0$ y todo punto $x'$ de $X'=X[\mathrm{T}_1,\ldots,\mathrm{T}_n]$ situado sobre $x$, el morfismo
\oldpage[IV-3]{211}
$f':X'\to Y'=Y[\mathrm{T}_1,\ldots,\mathrm{T}_n]$, deducido de $f$ por cambio de base, es abierto en el punto $x'$.
Teniendo en cuenta el lema \sref{IV.14.4.1.1}, \sref{IV.2.3.4} y \sref{IV.4.2.7}, quedamos reducidos a demostrar que $f$ es \emph{abierto} en el punto $x$; además, por {\hyperlink{ega4.c04.prop.14-1-2}{(14.1.2, (iii))}}, basta evidentemente demostrar que la restricción de $f$ a un subpreesquema cerrado de $X$ cuyo espacio subyacente sea $Z$ es abierta en el punto $x$, de modo que podemos limitarnos al caso en que $X=Z$ es irreducible.
Reemplazando $X$ por un entorno abierto $V$ de $x$ tal que $V\cap X_y$ sea irreducible, podemos suponer, en virtud de \sref{IV.13.3.1}, que el morfismo $f$ se factoriza como $X\xrightarrow{g}Y''=Y[\mathrm{T}_1,\ldots,\mathrm{T}_m]\to Y$, donde $g$ es cuasifinito, dominante y localmente de tipo finito.
Como el morfismo estructural $Y''\to Y$ es abierto \sref{IV.2.4.6}, quedamos reducidos a demostrar que $g:X\to Y''$ es abierto en el punto $x$.
Por \sref{IV.14.4.1.1}, $g(x)$ es un punto geométricamente unirramificado de $Y''$.
Quedamos así reducidos a demostrar el lema siguiente:

\begin{lemma}[14.4.1.3]
\label{IV.14.4.1.3}
Sean $X$, $Y$ dos preesquemas irreducibles, y $f:X\to Y$ un morfismo localmente cuasifinito y dominante.
Si $x\in X$ es tal que $y=f(x)$ es un punto unirramificado de $Y$, entonces $f$ es abierto en el punto $x$.
\end{lemma}

Es suficiente demostrar que $f(\Spec(\sh{O}_{X,x}))=\Spec(\sh{O}_{Y,y})$ \sref{IV.1.10.3}.
Por tanto, mediante el cambio de base $\Spec(\sh{O}_{Y,y})\to Y$, podemos limitarnos al caso $Y=\Spec(A)$, donde $A$ es un anillo local y $y$ es el punto cerrado de $Y$ (teniendo en cuenta \sref[I]{I.3.6.5} y \sref[$0_{\mathrm{I}}$]{0.2.1.8}); esto prueba que $X\times_Y\Spec(\sh{O}_{Y,y})$ es irreducible; reemplazando $f$ por $f_{\mathrm{red}}$, podemos suponer que $X$ e $Y$ son reducidos, y por tanto íntegros.
Reemplazando finalmente $X$ e $Y$ por entornos afines de $x$ y $y$, respectivamente, podemos suponer, por \sref{IV.8.12.9}, que el morfismo $f$ se factoriza como $X\xrightarrow{j}X_1\xrightarrow{g}Y$, donde $j$ es una \emph{inmersión abierta} y $g$ un morfismo \emph{finito} (evidentemente dominante); como $X$ y $X_1$ son afines, $j$ es afín, luego separado y cuasicompacto, y por consiguiente $j$ se factoriza como $X\xrightarrow{u}X_2\xrightarrow{h}X_1$, donde $X_2$ es la imagen cerrada de $X$ por $j$, $h$ la inyección canónica y $u$ una inmersión abierta \sref[I]{I.9.5.3}.
En otros términos, podemos suponer que $X_1$ es íntegro, o incluso que $X_1=\Spec(B)$, donde $B$ es una $A$-álgebra íntegra y finita que contiene a $A$, puesto que $g$ es dominante. Si $\mathfrak{n}$ es el ideal primo de $B$ correspondiente al punto $x$, la hipótesis de que $A$ es unirramificado implica, por \sref{IV.14.4.1.2}, que el morfismo $\Spec(B_\mathfrak{n})\to\Spec(A)$ es sobreyectivo, es decir, que $\Spec(\sh{O}_{X,x})\to\Spec(\sh{O}_{Y,y})$ es sobreyectivo.
Q.E.D.

\begin{corollary}[14.4.2]
\label{IV.14.4.2}
Sea $f:X\to Y$ un morfismo localmente de tipo finito, $y$ un punto geométricamente unirramificado de $Y$, y $\eta$ el punto genérico de la única componente irreducible $Y_0$ de $Y$ que contiene a $y$.
Las condiciones siguientes son equivalentes:
\begin{enumerate}[label=\emph{\alph*)}]
  \item $f$ es universalmente abierto en todos los puntos de $X_y$ (o, lo que equivale a lo mismo, en los puntos maximales de $X_y$ {\hyperlink{ega4.c04.subsec.14-3-3-1}{(14.3.3.1, (ii))}}).
  \item Para todo $x\in X_y$, existe una componente irreducible $Z$ de $X$ que contiene a $x$ y es equidimensional sobre $Y$ en el punto $x$ \sref{IV.13.2.2}.
  \item[\emph{b$'$)}] Para todo $x\in X_y$ y todo entorno abierto $U$ de $x$ en $X$, se tiene
    \[
      \dim(U\cap X_{\eta})\geq\dim_x(U\cap X_y).
    \]
  \item[\emph{b$''$)}] Para todo abierto $U$ de $X$, se tiene $\dim(U\cap X_{\eta})\geq\dim(U\cap X_y)$.
\end{enumerate}
\end{corollary}
\oldpage[IV-3]{212}
Si además $Y$ es localmente noetheriano, estas condiciones son también equivalentes a la siguiente:
\begin{enumerate}[label=\emph{\alph*)}]
  \item[\emph{c)}] $f$ es abierto en todos los puntos de $X_y$ (o, lo que equivale a lo mismo, en los puntos maximales de $X_y$ {\hyperlink{ega4.c04.subsec.14-3-3-1}{(14.3.3.1, (ii))}}).
\end{enumerate}
La equivalencia de \emph{a)} y \emph{c)} cuando $Y$ es localmente noetheriano resulta de \sref{IV.14.4.1}; las condiciones \emph{b)} o \emph{b$'$)}, aplicadas a los puntos \emph{maximales} de $X_y$, implican \emph{a)} también en virtud de {\hyperlink{ega4.c04.subsec.14-3-3-1}{(14.3.3.1, (ii))}}; finalmente, \emph{b$'$)} y \emph{b$''$)} son equivalentes, ya que
\[
  \dim(U\cap X_y) = \sup_{x\in X_y}(\dim_x(U\cap X_y)).
\]

Queda por ver que la condición \emph{a)} implica \emph{b)} y \emph{b$'$)} en \emph{todo} punto $x\in X_y$.
Sea $d=\dim_x(X_y)$, y sea $x'$ el punto genérico de una componente irreducible de $X_y$ que contiene a $x$ y tiene dimensión $d$.
En virtud de \emph{a)} y de \sref{IV.14.4.1}, existe una componente irreducible $Z$ de $X$ que contiene a $x'$ y es equidimensional sobre $Y$ en el punto $x'$, de modo que $\dim_{x'}(Z\cap X_y)=\dim(Z\cap X_\eta)$.
Pero, por construcción, $\dim_{x'}(Z\cap X_y)=\dim_x(X_y)=d$, y $\dim_x(Z\cap X_y)\leq\dim(Z\cap X_y)=d$; teniendo en cuenta \sref{IV.13.1.6}, esto prueba que $Z$ es equidimensional sobre $Y$ en el punto $x$; por tanto, \emph{a)} implica \emph{b)}.
Además, tenemos $\dim(X_\eta)\geq\dim(Z\cap X_\eta)=d=\dim_x(X_y)$.
Reemplazando $X$ por un entorno abierto $U$ de $x$, vemos así que \emph{a)} implica \emph{b$'$)}.
Q.E.D.

\begin{corollary}[14.4.3]
\label{IV.14.4.3}
Sea $Y$ un preesquema geométricamente unirramificado, y $f:X\to Y$ un morfismo localmente de tipo finito.
Las condiciones siguientes son equivalentes:
\begin{enumerate}[label=\emph{\alph*)}]
  \item $f$ es universalmente abierto.
  \item Para todo abierto $U$ de $X$, todo $y\in Y$ y toda generalización $y'$ de $y$, se tiene
    \[
      \dim(U\cap X_{y'})\geq\dim(U\cap X_y).
    \]
\end{enumerate}
Si además $Y$ es localmente noetheriano, estas condiciones son también equivalentes a la siguiente:
\begin{enumerate}[label=\emph{\alph*)}]
  \item[\emph{c)}] $f$ es abierto.
\end{enumerate}
\end{corollary}

\begin{corollary}[14.4.4]
\label{IV.14.4.4}
(Criterio de Chevalley).
Sea $f:X\to Y$ un morfismo localmente de tipo finito.
\begin{enumerate}[label=(\roman*)]
  \item Si $f$ es equidimensional en un punto $x\in X$ \sref{IV.13.3.2} y si $y=f(x)$ es un punto geométricamente unirramificado de $Y$, entonces $f$ es universalmente abierto en el punto $x$.
  \item Si $Y$ es geométricamente unirramificado, $f$ es universalmente abierto en todos los puntos de $X$ donde $f$ es equidimensional, y el conjunto de estos puntos es abierto en $X$.
 En particular, si $f$ es equidimensional, es universalmente abierto.
\end{enumerate}
\end{corollary}

La afirmación (ii) es una consecuencia trivial de (i), pues ya sabemos que el conjunto de puntos donde $f$ es equidimensional es abierto \sref{IV.13.3.2}.
En cuanto a la afirmación (i), resulta del hecho de que la hipótesis implica que la condición \emph{b)} de \sref{IV.14.4.1} se satisface en el punto genérico de una componente irreducible de $X_y$ que contiene a $x$, teniendo en cuenta \sref{IV.13.3.1}; basta, por tanto, aplicar \sref{IV.14.4.1}.

\begin{remark}[14.4.5]
\label{IV.14.4.5}
Se puede demostrar que, si $Y$ es localmente noetheriano y todas las generalizaciones de $y=f(x)$ son puntos geométricamente unirramificados de $Y$ (cf.\ \sref{IV.6.15.2}), entonces, si $f$ es equidimensional en el punto $x$, es universalmente abierto \emph{en un entorno de $x$}.
\end{remark}

\begin{corollary}[14.4.6]
\label{IV.14.4.6}
\oldpage[IV-3]{213}
Sea $Y$ un preesquema localmente noetheriano, y $f:X\to Y$ un morfismo localmente de tipo finito.
Sea $y$ un punto geométricamente unirramificado de $Y$, y supongamos además que, para todo $x\in f^{-1}(y)$, el anillo $\sh{O}_{X,x}$ es equidimensional.
Entonces las condiciones siguientes son equivalentes:
\begin{enumerate}[label=\emph{\alph*)}]
  \item $f$ es equidimensional \sref{IV.13.3.2} en todos los puntos de $f^{-1}(y)$.
  \item $f$ es abierto en todos los puntos de $f^{-1}(y)$.
  \item $f$ es universalmente abierto en todos los puntos de $f^{-1}(y)$.
\end{enumerate}
\end{corollary}

En efecto, \emph{c)} implica trivialmente \emph{b)}, y \emph{a)} implica \emph{c)} en virtud de \sref{IV.14.4.4}; finalmente, debido a las hipótesis sobre $\sh{O}_y$ y $\sh{O}_x$, \emph{b)} implica \emph{a)} por \sref{IV.14.2.2}.
Más generalmente:

\begin{proposition}[14.4.7]
\label{IV.14.4.7}
Sea $Y$ un preesquema localmente noetheriano, $f:X\to Y$ un morfismo localmente de tipo finito, y $x$ un punto de $X$, tal que $y=f(x)$ es geométricamente unirramificado.
Las condiciones siguientes son equivalentes:
\begin{enumerate}[label=\emph{\alph*)}]
  \item $f$ es equidimensional \sref{IV.13.3.2} en el punto $x$.
  \item El anillo $\sh{O}_x$ es equidimensional y $f$ es abierto en los puntos genéricos de las componentes irreducibles de $f^{-1}(y)$ que contienen a $x$ (y, por tanto, también en todo punto de una componente de ese tipo).
  \item El anillo $\sh{O}_x$ es equidimensional y $f$ es universalmente abierto en los puntos genéricos de las componentes irreducibles de $f^{-1}(y)$ que contienen a $x$ (y, por tanto, también en todo punto de una componente de ese tipo).
\end{enumerate}
Además, cuando se satisfacen estas condiciones, para todo subpreesquema cerrado reducido $X_i$ de $X$ cuyo espacio subyacente sea una componente irreducible de $X$ que contiene a $x$, la restricción $f_i:X_i\to Y$ de $f$ es un morfismo equidimensional en el punto $x$, y universalmente abierto en todos los puntos de las componentes irreducibles de $f^{-1}(y)\cap X_i$ que contienen a $x$.
\end{proposition}

La condición \emph{a)} implica que $f$ es equidimensional en todos los puntos genéricos de las componentes irreducibles de $f^{-1}(y)$ que contienen a $x$ \sref{IV.13.3.1}, y, por consiguiente, en virtud de \sref{IV.14.4.4}, universalmente abierto en esos puntos; el mismo argumento aplicado a cada $f_i$ (teniendo en cuenta \sref{IV.13.3.3}) demuestra la última afirmación de la proposición, a la vista de {\hyperlink{ega4.c04.subsec.14-3-3-1}{(14.3.3.1, (iii))}}.
Además, por \sref{IV.14.2.1}, tenemos las relaciones
\begin{align}
  \dim(\sh{O}_{X_i,x}) &= \dim(\sh{O}_y) + \dim(\sh{O}_{X_i,x}\otimes_{\sh{O}_y}\kappa(y))
  \tag{14.4.7.1}\label{IV.14.4.7.1}\\
  \dim(\sh{O}_{X,x}) &= \dim(\sh{O}_y) + \dim(\sh{O}_{X,x}\otimes_{\sh{O}_y}\kappa(y))
  \tag{14.4.7.2}\label{IV.14.4.7.2}
\end{align}
Y como $f$ es equidimensional en el punto $x$, se deduce de \sref{IV.13.3.1} que $\dim_x(X_i\cap f^{-1}(y))=\dim_x(f^{-1}(y))$, y por tanto, mediante \sref{IV.5.2.3}, $\dim(\sh{O}_{X_i,x}\otimes_{\sh{O}_y}\kappa(y))=\dim(\sh{O}_{X,x}\otimes_{\sh{O}_y}\kappa(y))$.
Concluimos que $\dim(\sh{O}_{X_i,x})=\dim(\sh{O}_{X,x})$ para todo $i$; dicho de otro modo, $\sh{O}_{X,x}$ es equidimensional, y esto termina de demostrar que \emph{a)} implica \emph{c)}.
Es claro que \emph{c)} implica \emph{b)}; finalmente, \emph{b)} implica la relación \eqref{IV.14.4.7.2} en virtud de \sref{IV.14.2.1}; se deduce entonces de \sref{IV.13.3.6} que \emph{b)} implica \emph{a)}.

\begin{proposition}[14.4.8]
\label{IV.14.4.8}
Sea $Y$ un preesquema noetheriano, $f:X\to Y$ un morfismo localmente de tipo finito, $y$ un punto de $Y$, y $x$ un punto maximal de $X_y=f^{-1}(y)$.
Las condiciones siguientes son equivalentes:
\begin{enumerate}[label=\emph{\alph*)}]
  \item El morfismo $f$ es universalmente abierto en $x$; dicho de otro modo, para todo cambio de base $g:Y'\to Y$ se satisface la propiedad $\mathbf{P}(Y')$: para todo punto $x'$ de $X'=X\times_Y Y'$ situado sobre $x$, el morfismo $f'=f_{(Y')}:X'\to Y'$ es abierto en $x'$.
\oldpage[IV-3]{214}
  \item[\emph{a$'$)}] La propiedad $\mathbf{P}(Y')$ se satisface para todo morfismo finito $g:Y'\to Y$.
  \item[\emph{a$''$)}] La propiedad $\mathbf{P}(Y'')$ se satisface para la normalización $Y''$ de $Y_{\mathrm{red}}$ \sref[II]{II.6.3.8}.
  \item Para todo punto $x''$ de $X''=X\times_Y Y''$ situado sobre $x$, existe una componente irreducible $Z''$ de $X''$ que contiene a $x''$ y es equidimensional sobre $Y''$ en el punto $x''$.
\end{enumerate}
\end{proposition}

Es trivial que \emph{a)} implica \emph{a$'$)}.
Para demostrar que \emph{a$'$)} implica \emph{a$''$)}, observemos que podemos escribir $Y''=\Spec(\sh{B})$, donde $\sh{B}$ es una $\sh{O}_Y$-álgebra cuasicoherente entera sobre $\sh{O}_Y$; como $Y$ es noetheriano, $\sh{B}$ es el límite inductivo de sus sub-$\sh{O}_Y$-álgebras $\sh{B}_\lambda$, que son cuasicoherentes y de tipo finito \sref[I]{I.9.6.6}; pero entonces las $\sh{B}_\lambda$ son $\sh{O}_Y$-álgebras \emph{finitas} \sref[II]{II.6.1.2}; por tanto podemos escribir $Y''=\varprojlim Y'_\lambda$, donde $Y'_\lambda=\Spec(\sh{B}_\lambda)$, y en consecuencia $X''=X\times_Y Y''=\varprojlim X'_\lambda$, con $X'_\lambda=X\times_Y Y'_\lambda$.
En virtud de \emph{a$'$)}, los morfismos $f'_\lambda:X'_\lambda\to Y'_\lambda$ son abiertos en todos los puntos de $X'_\lambda$ situados sobre $x$; concluimos que $f''$ es abierto en todos los puntos de $X''$ situados sobre $x$, por \sref{IV.8.10.1} y {\hyperlink{ega4.c04.subsec.14-3-3-1}{(14.3.3.1, (i))}}.

Como el preesquema $Y''$ es \emph{normal} por definición, el hecho de que \emph{b)} implique \emph{a$''$)} se deduce de \sref{IV.14.4.4}, aplicado a la componente irreducible equidimensional del enunciado y a la restricción de $f''$ a dicha componente.
Queda por demostrar que \emph{a$''$)} implica \emph{a)} y \emph{b)}.
Teniendo en cuenta \sref{IV.1.10.3}, podemos limitarnos al caso $Y=\Spec(\sh{O}_y)$, observando que el morfismo canónico $\Spec(\sh{O}_y)\to Y$ es universalmente bicontinuo \sref[I]{I.3.6.5}, mientras que $Y''\times_Y\Spec(\sh{O}_y)$ es la normalización de $(\Spec(\sh{O}_y))_{\mathrm{red}}$, como se deduce de la conmutatividad de las operaciones de cierre entero y localización (Bourbaki, \emph{Alg.\ comm.}, cap.~V, \S~1, n\textsuperscript{o}~5, prop.~16).
Supongamos, pues, que $Y=\Spec(A)$, donde $A$ es un anillo local noetheriano, y $y$ es el punto cerrado de $Y$; sabemos por \sref[0]{0.23.2.5} que existe una factorización
\[
  Y'' \xrightarrow{u} Y_1 \xrightarrow{v} Y
\]
del morfismo estructural, tal que $v$ es un morfismo \emph{finito} y sobreyectivo, mientras que $u$ es entero, \emph{radicial} y dominante, luego sobreyectivo \sref[II]{II.6.1.10} y, por consiguiente, un \emph{homeomorfismo} universal \sref{IV.2.4.5}.
Si ponemos $X_1=X\times_Y Y_1$, la proyección $X''\to X_1$ es, por tanto, un homeomorfismo, y la hipótesis \emph{a$''$)} implica que $f_1=f_{(Y_1)}:X_1\to Y_1$ es abierto en todos los puntos de $X_1$ situados sobre $x$.
Además, esto muestra que, para demostrar la propiedad \emph{b)}, basta demostrar la misma propiedad reemplazando $Y''$, $X''$ y $x''$ por $Y_1$, $X_1$ y un punto $x_1$ de $X_1$ situado sobre $x$.
Pero $Y_1$ es noetheriano y, además, geométricamente unirramificado, puesto que $Y''$ es normal y $u$ es radicial \sref{IV.6.15.1}; la propiedad \emph{b)} resulta, por tanto, de \sref{IV.14.4.1}.
Queda por demostrar que $f$ es universalmente abierto en el punto $x$, lo que resultará del lema siguiente:

\begin{lemma}[14.4.8.1]
\label{IV.14.4.8.1}
Sea $v:Y_1\to Y$ un morfismo cerrado (resp.\ universalmente cerrado) y sobreyectivo.
Para que un morfismo $f:X\to Y$ sea abierto (resp.\ universalmente abierto) en un punto $x\in X$, basta que $f_1=f_{(Y_1)}:X_1=X\times_Y Y_1\to Y_1$ sea abierto (resp.\ universalmente abierto) en todos los puntos de $X_1$ situados sobre $x$.
\end{lemma}

La segunda afirmación resulta trivialmente de la primera y del hecho de que, para todo cambio de base $Y'\to Y$, el morfismo $v_{(Y')}:Y_1\times_Y Y'\to Y'$ sigue siendo sobreyectivo y universalmente cerrado.
Para demostrar la primera afirmación, consideremos
\oldpage[IV-3]{215}
un entorno abierto $U$ de $x$ en $X$; como $v$ es cerrado y sobreyectivo, para que $f(U)$ sea un entorno de $y=f(x)$ es necesario y suficiente que $v^{-1}(f(U))$ sea un entorno de $v^{-1}(y)$.
Pero si $p:X_1\to X$ es la proyección canónica, tenemos $v^{-1}(f(U))=f_1(p^{-1}(U))$ \sref[I]{I.3.4.8}, y la hipótesis implica que $f_1(p^{-1}(U))$ es un entorno de $f_1(p^{-1}(x))=v^{-1}(y)$ \sref[I]{I.3.4.8}.

\begin{corollary}[14.4.9]
\label{IV.14.4.9}
Sea $Y$ un preesquema noetheriano, y $f:X\to Y$ un morfismo localmente de tipo finito.
Las condiciones siguientes son equivalentes:
\begin{enumerate}[label=\emph{\alph*)}]
  \item $f$ es universalmente abierto; dicho de otro modo, para todo cambio de base $Y'\to Y$, el morfismo $f_{(Y')}:X\times_Y Y'\to Y'$ es abierto.
  \item[\emph{a$'$)}] Para todo morfismo finito $Y_1\to Y$, $f_{(Y_1)}$ es abierto.
  \item[\emph{a$''$)}] Si $Y''$ es la normalización de $Y_{\mathrm{red}}$, $f_{(Y'')}$ es abierto.
  \item Para todo punto $x''$ de $X''=X\times_Y Y''$, existe una componente irreducible $Z''$ de $X''$ que contiene a $x''$ y es equidimensional sobre $Y''$ en el punto $x''$ (cf.\ {\hyperlink{ega4.c04.item.14-4-10}{(14.4.10, (ii))}}).
\end{enumerate}
\end{corollary}

Esto resulta inmediatamente de \sref{IV.14.4.8} y \sref{IV.14.1.4}.

\begin{remarks}[14.4.10]
\label{IV.14.4.10}
(i) La equivalencia de las condiciones \emph{a)} y \emph{b)} en \sref{IV.14.4.8} (resp.\ \sref{IV.14.4.9}) es también válida para todo preesquema $Y$ y todo morfismo $f$ localmente de tipo finito.
En efecto, \emph{a)} implica \emph{b)} en virtud de \sref{IV.14.4.1}; recíprocamente, \emph{b)} implica que $f''$ es universalmente abierto en los puntos de $X''$ situados sobre $x$ (resp.\ en todo punto de $X''$) en virtud de \sref{IV.14.4.1}, y concluimos la propiedad \emph{a)} aplicando el lema \sref{IV.14.4.8.1} al morfismo entero y sobreyectivo $Y''\to Y$.

Puede que, en \sref{IV.14.4.1}, la hipótesis suplementaria de que $Y$ es noetheriano sea superflua para la equivalencia de \emph{a)} y \emph{c)} (cf.\ \sref{IV.14.3.14}).
En tal caso, las hipótesis de noetherianidad también serían superfluas en \sref{IV.14.4.2}, \sref{IV.14.4.3}, \sref{IV.14.4.8} y \sref{IV.14.4.9}.

(ii) Se pueden dar ejemplos de morfismos $f:X\to Y$ con las propiedades siguientes: $Y$ es noetheriano, regular y de dimensión~2, $f$ es \emph{universalmente abierto} y de tipo finito, $X$ tiene dos componentes irreducibles $X_1$, $X_2$, pero la restricción $X_1\to Y$ de $f$ a una de ellas \emph{no es un morfismo abierto}.
El principio de la construcción se basa en el método general de ``pegado'', que se presentará en el cap.~V y que aquí solo podemos esbozar.
Comencemos con un punto cerrado $y$ de $Y$, y consideremos el $Y$-esquema $Y_1$ obtenido haciendo la explosión de $y$ \sref[II]{II.8.1.3}. Si $f_1:Y_1\to Y$ es el morfismo estructural, sabemos que la restricción de $f_1$ a $f_1^{-1}(Y-\{y\})$ es un \emph{isomorfismo} sobre $Y-\{y\}$ (\emph{loc.\ cit.}), mientras que la fibra $f_1^{-1}(y)$ es isomorfa a $\operatorname{Proj}(S)$, donde $S=\bigoplus_{k=0}^{\infty}\mathfrak{m}_y^k/\mathfrak{m}_y^{k+1}$ \sref[II]{II.3.5.3}, es decir, aquí $\mathbf{P}^1_{\kappa(y)}$; se deduce de \sref{IV.14.4.1} que $f_1$ no es abierto en el punto genérico de $f_1^{-1}(y)$.
Por otra parte, pongamos $Y_2=\mathbf{P}^1_Y$, y sea $f_2:Y_2\to Y$ el morfismo estructural. Se deduce de \sref[II]{II.8.4.4} que $f_2$ es plano y, por tanto, universalmente abierto \sref{IV.2.4.6}; además, por \sref[II]{II.3.5.3}, $f_2^{-1}(y)$ es isomorfo a $\mathbf{P}^1_{\kappa(y)}$. Basta entonces ``pegar'' $Y_1$ y $Y_2$ a lo largo de las fibras isomorfas $f_1^{-1}(y)$ y $f_2^{-1}(y)$; esto da un morfismo $f:X\to Y$ cuyas componentes irreducibles $X_1$, $X_2$ se identifican canónicamente con $Y_1$, $Y_2$, respectivamente, y cuyas restricciones a estas componentes son $f_1$, $f_2$.

Recordemos, sin embargo, \sref{IV.12.1.1.5} que, si $Y$ es localmente noetheriano, $f:X\to Y$ es de tipo finito y \emph{plano}, toda componente irreducible de $X$ es equidimensional sobre $Y$ (y, en consecuencia, la restricción de $f$ a una componente de ese tipo es \emph{universalmente abierta} si todos los puntos de $Y$ son \emph{geométricamente unirramificados}).

(iii) Recordemos {\hyperlink{ega4.c04.item.12-1-2}{(12.1.2, (i))}} que existen morfismos $f:X\to Y$ con las propiedades siguientes: $Y$ es noetheriano (no geométricamente unirramificado), $f$ es finito y plano (e incluso \'etale \sref{IV.17.6.3}), pero la restricción de $f$ a una componente irreducible de $X$ \emph{no es un morfismo abierto} (aunque $f$ mismo sí lo es por \sref{IV.2.4.6}).

(iv) El criterio de Chevalley \sref{IV.14.4.4} explica la importancia de la noción de morfismo universalmente abierto.
Esta noción permite, en efecto, sustituir en numerosos resultados, más o menos clásicos, una hipótesis de \emph{normalidad} por la hipótesis de que cierto morfismo es universalmente abierto; el enunciado más general
\oldpage[IV-3]{216}
así obtenido se aplicará en particular a los morfismos \emph{planos} \sref{IV.2.4.6}, cuya importancia en geometría algebraica va en aumento.
Puede considerarse que los enunciados que contienen la hipótesis de que un morfismo es universalmente abierto son generalizaciones comunes de enunciados que contienen una hipótesis de normalidad y de enunciados que contienen una hipótesis de planitud.
\end{remarks}


\subsection{Morfismos universalmente abiertos y cuasisecciones}
\label{subsection:IV.14.5}
\label{IV.14.5}

\begin{lemma}[14.5.1]
\label{IV.14.5.1}
Sea $Y$ un preesquema irreducible y localmente noetheriano, $f:X\to Y$ un morfismo localmente de tipo finito, y $x$ un punto de $X$, tales que $f$ es equidimensional \sref{IV.13.3.2} en el punto $x$; pongamos $y=f(x)$, $e=\dim_x(f^{-1}(y))$.
Sea $X'$ un subespacio cerrado irreducible de $X$ que contiene a $x$, y sea $n$ un entero tal que $\codim(X',X)\leq n$ y $\dim_x(X'\cap f^{-1}(y))\leq e-n$.
Entonces necesariamente $\codim(X',X)=n$, $\dim_x(X'\cap f^{-1}(y))=e-n$, y la restricción $X'\to Y$ de $f$ es un morfismo equidimensional en $x$ (y, a fortiori, dominante).
\end{lemma}

Como la cuestión es local en $X$, podemos suponer que $f$ es de tipo finito y \emph{equidimensional}, de modo que para \emph{todo} $z\in Y$, \emph{todas} las componentes irreducibles de $f^{-1}(z)$ tienen dimensión $e$ {\hyperlink{ega4.c04.prop.13-3-1}{(13.3.1, $a''$))}}.
Sea $x'$ el punto genérico de $X'$, $y'=f(x')$, $Y'=\overline{\{y'\}}=\overline{f(X')}$ y pongamos $Z=f^{-1}(Y')$.
En virtud de \sref[0]{0.14.2.2}, tenemos
\begin{equation}
  \codim(X',Z)\leq n
  \tag{14.5.1.1}\label{IV.14.5.1.1}
\end{equation}
 y, si los dos miembros son iguales, necesariamente $\codim(X',X)=n$ y $Z$ contiene una componente irreducible de $X$, lo que implica por \sref{IV.13.3.1} que $f(X')$ es denso en $Y$ y, por tanto, $Z=X$; así, la igualdad $\codim(X',Z)=n$ equivale a la conjunción de la igualdad $\codim(X',X)=n$ y la relación $Z=X$.

Por otra parte, razonando en los preesquemas reducidos de $Y$ y $X$ que tienen respectivamente a $Y'$ y $Z$ como espacios subyacentes, deducimos de \sref{IV.5.1.2} y de \sref[I]{I.3.6.5} que
\begin{equation}
  \codim(X'\cap f^{-1}(y'),f^{-1}(y'))=\codim(X',Z).
  \tag{14.5.1.2}\label{IV.14.5.1.2}
\end{equation}
Por hipótesis, $f^{-1}(y')$ es biequidimensional \sref{IV.5.2.1} y de dimensión $e$, luego, por \sref[0]{0.14.3.5}, en virtud de \eqref{IV.14.5.1.2} y \eqref{IV.14.5.1.1},
\begin{equation}
  e'=\dim(X'\cap f^{-1}(y'))=e-\codim(X',Z)\geq e-n
  \tag{14.5.1.3}\label{IV.14.5.1.3}
\end{equation}
la igualdad se da si y solo si $\codim(X',X)=n$ y $f(X')$ es denso en $Y$.
Finalmente, por \sref{IV.13.1.1}, tenemos $\dim_x(X'\cap f^{-1}(y))\geq e'$, de donde, por \eqref{IV.14.5.1.3},
\begin{equation}
  \dim_x(X'\cap f^{-1}(y))\geq e'\geq e-n.
  \tag{14.5.1.4}\label{IV.14.5.1.4}
\end{equation}
Pero, por hipótesis, también tenemos $\dim_x(X'\cap f^{-1}(y))\leq e-n$, de donde se siguen las conclusiones de la proposición.

\begin{corollary}[14.5.2]
\label{IV.14.5.2}
Bajo las hipótesis de \sref{IV.14.5.1} sobre $f$, $X$, $Y$ y $x$, supongamos además que $x$ no es un punto maximal de $f^{-1}(y)$.
Entonces existen un entorno abierto afín $U$ de $x$ en $X$ y una sección $g\in\Gamma(U,\sh{O}_X)$ tales que el conjunto $X'$ de los $x'\in U$ que satisfacen $g(x')=0$ contiene a $x$ y no contiene ningún punto maximal de $f^{-1}(y)$.
Para todo $g$ con estas propiedades, $X'$ es equidimensional sobre $Y$ en el punto $x$, y tenemos
\begin{equation}
  \dim_x(X'\cap f^{-1}(y))=e-1 \qquad\text{y}\qquad \codim(X',X)=1.
  \tag{14.5.2.1}\label{IV.14.5.2.1}
\end{equation}
\end{corollary}

\oldpage[IV-3]{217}
Podemos limitarnos al caso en que $X=U$ es un entorno abierto afín de $x$ tal que todas las componentes irreducibles de $f^{-1}(y)$ contienen a $x$.
Estas componentes corresponden a los ideales primos mínimos de $\sh{O}_x/\mathfrak{m}_y\sh{O}_x$, y por hipótesis estos ideales son distintos de $\mathfrak{m}_x/\mathfrak{m}_y\sh{O}_x$ (Bourbaki, \emph{Alg.\ comm.}, cap.~II, \S~1, n\textsuperscript{o}~1, prop.~2); para obtener un $g\in\Gamma(U,\sh{O}_X)$ que satisfaga las condiciones del enunciado, basta tomar $g\in\mathfrak{j}_x$ tal que la imagen de $g$ en $\mathfrak{m}_x$ no pertenezca a ninguno de los ideales primos precedentes.
Además, como $\codim(X',X)\leq 1$ \sref{IV.5.1.8}, y como $X'$ no contiene ninguna de las componentes irreducibles de $f^{-1}(y)$ de dimensión $e$, y puesto que estas son todas las componentes irreducibles de $f^{-1}(y)$, tenemos, por \sref[0]{0.14.2.2.2}, $\dim_x(X'\cap f^{-1}(y))\leq e-1$.
Basta entonces aplicar \sref{IV.14.5.1}.

\begin{proposition}[14.5.3]
\label{IV.14.5.3}
Sea $Y$ un preesquema irreducible y localmente noetheriano, $f:X\to Y$ un morfismo localmente de tipo finito, y $x$ un punto de $X$.
Supongamos que $f$ es equidimensional en el punto $x$ y que $x$ es cerrado en $f^{-1}(f(x))$.
Entonces existe un subconjunto irreducible $X'$ de $X$, localmente cerrado en $X$, que contiene a $x$, tal que la restricción $X'\to Y$ de $f$ (donde $X'$ está dotado de la estructura de subesquema reducido inducida por $X$) es un morfismo \emph{cuasifinito y dominante}.
\end{proposition}

En efecto, con la notación de \sref{IV.14.5.2}, la hipótesis de que $x$ es cerrado en $f^{-1}(y)$ implica que $x$ no es un punto maximal de $X'\cap f^{-1}(y)$ mientras que $e-1\geq 1$.
Basta, por tanto, aplicar \sref{IV.14.5.2} por inducción descendente sobre $e=\dim_x(f^{-1}(f(x)))$ hasta llegar a $e=1$; la aplicación de \sref{IV.14.5.2} en este último caso proporciona un subespacio $X'$ tal que $X'\cap f^{-1}(f(x))$ es noetheriano y de dimensión $0$, luego finito y discreto; como $X'$ es equidimensional sobre $Y$, $X'\cap f^{-1}(f(x'))$ tiene dimensión $0$ para todo $x'\in X'$, lo que implica que la restricción $X'\to Y$ de $f$ es un morfismo cuasifinito \sref[II]{II.6.2.2}.

\begin{corollary}[14.5.4]
\label{IV.14.5.4}
Bajo las hipótesis de \sref{IV.14.5.3}, supongamos además que $Y=\Spec(A)$, donde $A$ es un anillo local noetheriano, íntegro y completo, y que $y=f(x)$ es el único punto cerrado de $Y$.
Entonces existe un anillo local íntegro $A'$, que contiene a $A$ y es una $A$-álgebra finita, tal que, si ponemos $Y'=\Spec(A')$ y $X'=X\times_Y Y'$, existe una $Y'$-sección $h:Y'\to X'$ cuyo morfismo compuesto $Y'\xrightarrow{h}X'\to X$ es una inmersión cuya imagen contiene a $x$.
\end{corollary}

Reemplazando si es necesario $X$ por un subpreesquema cerrado irreducible reducido de $X$, podemos, en virtud de \sref{IV.14.5.3}, limitarnos al caso en que el morfismo $f$ ya es cuasifinito y dominante.
Usando \sref[II]{II.6.2.5}, deducimos que $\sh{O}_x=A'$ es un anillo íntegro que es una $A$-álgebra finita, y que $X$ es una suma disjunta del subpreesquema cerrado $Y'=\Spec(A')$ y de un subpreesquema $Y''$; el esquema $Y'$ satisface lo pedido, pues el morfismo compuesto $Y'\to X'\to X$ no es otro que el morfismo canónico $\Spec(\sh{O}_x)\to X$.

\begin{remark}[14.5.5]
\label{IV.14.5.5}
Si, en el enunciado de \sref{IV.14.5.4}, no se exige que el morfismo $Y'\to X$ sea una inmersión, se puede suponer además que $A'$ es \emph{íntegramente cerrado}: en efecto, basta reemplazar $A'$ por su \emph{cierre entero} $A_1$, pues sabemos por \sref[0]{0.23.1.5} que $A_1$ es un $A$-módulo de tipo finito.
\end{remark}

\begin{proposition}[14.5.6]
\label{IV.14.5.6}
Sea $Y$ un preesquema localmente noetheriano, irreducible, regular y de dimensión $1$, $f:X\to Y$ un morfismo localmente de tipo finito, y $y$ un punto de $Y$.
Las condiciones siguientes son equivalentes:
\begin{enumerate}[label=\emph{\alph*)}]
  \item $f_{\mathrm{red}}$ es plano en todo punto de $f^{-1}(y)$.
\oldpage[IV-3]{218}
  \item $f$ es universalmente abierto en un entorno de $f^{-1}(y)$.
  \item $f$ es abierto en un entorno de $f^{-1}(y)$.
  \item Toda componente irreducible de $X$ que corta a $f^{-1}(y)$ domina a $Y$.
  \item Para todo punto $x\in f^{-1}(y)$ cerrado en $f^{-1}(y)$, y toda componente irreducible $X_i$ de $X$ que contiene a $x$, existe un subconjunto irreducible $X'$ de $X$, localmente cerrado en $X$, que contiene a $x$, está contenido en $X_i$, y tal que la restricción $X'\to Y$ de $f$ es un morfismo cuasifinito y dominante.
\end{enumerate}
\end{proposition}

La equivalencia de \emph{a)}, \emph{b)}, \emph{c)} y \emph{d)} ya se ha demostrado en \sref{IV.14.3.8}.
Es claro que \emph{e)} implica \emph{d)}, pues para toda componente irreducible $X_i$ de $X$ que corta a $f^{-1}(y)$ existe en $X_i\cap f^{-1}(y)$ un punto cerrado \sref{IV.5.1.11}.
Finalmente, para demostrar que \emph{c)} implica \emph{e)}, podemos limitarnos al caso en que $f$ es un morfismo de tipo finito; consideremos un punto cerrado $x$ de $f^{-1}(y)$ y sea $X_i$ una componente irreducible de $X$ que contiene a $x$.
Como la restricción $f_i:X_i\to Y$ de $f$ a $X_i$ es un morfismo dominante, de la equivalencia de \emph{c)} y \emph{d)} para $f_i$ se deduce que este morfismo es abierto en los puntos genéricos de $X_i\cap f^{-1}(y)$. Se deduce de \sref{IV.14.2.2} que $f_i$ es equidimensional en el punto $x$, y concluimos con ayuda de \sref{IV.14.5.3}.

\begin{remark}[14.5.7]
\label{IV.14.5.7}
Si, en el enunciado de \sref{IV.14.5.5}, suponemos que $Y=\Spec(A)$, donde $A$ es un anillo de valoración discreta completo, y que $y$ es el punto cerrado de $Y$, podemos suponer además que $X'=\Spec(A')$, donde $A'$ es un anillo de valoración discreta que es una $A$-álgebra finita, como se deduce de la demostración de \sref{IV.14.5.4} y del hecho de que un anillo local íntegro, regular y de dimensión $1$ es un anillo de valoración discreta \sref[II]{II.7.1.6}.
\end{remark}

\begin{proposition}[14.5.8]
\label{IV.14.5.8}
Sea $Y$ un preesquema localmente noetheriano, $f:X\to Y$ un morfismo localmente de tipo finito, y $y$ un punto de $Y$.
Para todo $Y$-preesquema $Y'$, pongamos $X(Y')=\Hom_Y(Y',X)$.
Sea $x$ un punto de $f^{-1}(y)$; para que $f$ sea universalmente abierto en $x$, es necesario y suficiente que se satisfaga la condición siguiente:

Para todo anillo de valoración discreta completo $A$ con cuerpo residual algebraicamente cerrado, todo morfismo $g:Y'=\Spec(A)\to Y$ tal que la imagen por $g$ del punto cerrado $y'$ de $Y'$ sea igual a $y$, y todo elemento $u_0\in X(\Spec(\kappa(y')))$ tal que $u_0(y')=x$, existen un anillo de valoración discreta $B$, un homomorfismo local $A\to B$ que hace de $B$ una $A$-álgebra finita y, si ponemos $Z'=\Spec(B)$, un elemento $u\in X(Z')$ tal que, si $z'$ es el punto cerrado de $Z'$, el diagrama
\[
  \xymatrix{
    \Spec(\kappa(z')) \ar[r] \ar[d]^{\ell} & Z' \ar[d]^{u} \\
    \Spec(\kappa(y')) \ar[r]_-{u_0} & X
  }
\]
es conmutativo.
\end{proposition}

\oldpage[IV-3]{219}
Observemos que, si $A$ y $B$ satisfacen las condiciones del enunciado, $B$ es un anillo de valoración discreta completo (Bourbaki, \emph{Alg.\ comm.}, cap.~III, \S~3, n\textsuperscript{o}~3, prop.~7 y cap.~IV, \S~2, n\textsuperscript{o}~2, cor.~3 de la prop.~9) cuyo cuerpo residual es isomorfo al de $A$ y, por tanto, algebraicamente cerrado. Además, como $u(z')=x$, existe en $X''=X\times_Y Z'$ un punto $x''$ cuyas proyecciones en $X$ y $Z'$ son, respectivamente, $x$ y $z'$, y que es racional sobre $\kappa(z')$; asimismo, como existe una $Z'$-sección $v$ de $X''$ tal que $v(z')=x''$, la imagen por $v$ del punto genérico $s$ de $Z'$ es una generalización $t$ de $x''$ cuya proyección en $Z'$ es $s$. Aplicando \sref{IV.14.3.6}, vemos que la condición del enunciado es \emph{suficiente}.
Demostremos ahora que es \emph{necesaria}.
Pongamos $X'=X\times_Y Y'$, $f'=f_{(Y')}:X'\to Y'$.
Por hipótesis existe un punto $x'\in X'$ situado sobre $x$ y $y'$ y racional sobre $\kappa(y')$ \sref[I]{I.3.3.14}, luego cerrado en $f'^{-1}(y')$.
En virtud de \sref{IV.14.3.6}, existe una componente irreducible $T'$ de $X'$ que contiene a $x'$ y domina a $Y'$; por tanto, por \sref{IV.14.3.8} y \sref{IV.14.3.13}, la restricción $T'\to Y'$ de $f'$ es equidimensional en $x'$.
Se deduce de \sref{IV.14.5.5} que existe una $A$-álgebra finita $B$ que es un dominio local íntegramente cerrado que domina a $A$ (y, por tanto, es un anillo de valoración discreta) y, poniendo $Z'=\Spec(B)$ y $X''=X\times_Y Z'=X'\times_{Y'}Z'$, una $Z'$-sección $v:Z'\to X''$ tal que la imagen del punto cerrado $z'$ de $Z'$ por el morfismo compuesto $Z'\xrightarrow{v}X''\to X'$ es $x'$; el morfismo compuesto $u:Z'\xrightarrow{v}X''\to X$ responde, por tanto, a la cuestión.

\begin{proposition}[14.5.9]
\label{IV.14.5.9}
Sea $Y$ un preesquema irreducible localmente noetheriano, $f:X\to Y$ un morfismo localmente de tipo finito, y sea $y$ un punto geométricamente unirramificado de $Y$.
Las condiciones siguientes son equivalentes:
\begin{enumerate}[label=\emph{\alph*)}]
  \item $f$ es universalmente abierto en todo punto de $f^{-1}(y)$.
  \item $f$ es abierto en todo punto de $f^{-1}(y)$.
  \item Para toda componente irreducible $Z$ de $f^{-1}(y)$, de punto genérico $z$, existe una componente irreducible $X_i$ de $X$ que contiene a $z$ y es equidimensional sobre $Y$ en $z$.
  \item Para todo punto cerrado $x$ de $f^{-1}(y)$, existe un subconjunto irreducible $X'$ de $X$, localmente cerrado en $X$, que contiene a $x$, y tal que la restricción $X'\to Y$ de $f$ es un morfismo cuasifinito dominante.
\end{enumerate}
\end{proposition}

La equivalencia de \emph{a)}, \emph{b)} y \emph{c)} ya se ha demostrado en \sref{IV.14.4.2}.
Para demostrar que \emph{a)} implica \emph{d)}, observemos que, por \sref{IV.14.4.2}, \emph{a)} implica que existe una componente irreducible $Z$ de $X$ que contiene a $x$ y es equidimensional sobre $Y$ en el punto $x$; la existencia de $X'$ resulta entonces de \sref{IV.14.5.3}, aplicado a la restricción $Z\to Y$ de $f$.
Recíprocamente, supongamos que se satisface \emph{d)}; por el criterio de Chevalley \sref{IV.14.4.4}, la restricción $X'\to Y$ de $f$ es universalmente abierta en $x$ y, \emph{a fortiori}, $f$ es abierto en $x$. Por tanto, $f$ es abierto en todo punto cerrado de $X_y=f^{-1}(y)$.
Pero $X_y$ es un $\kappa(y)$-preesquema localmente de tipo finito y, por tanto, un preesquema de Jacobson; el conjunto de puntos cerrados de $X_y$ es, pues, \emph{denso} en $X_y$ \sref{IV.10.3.1}, y se sigue de \sref{IV.14.1.4} que $f$ es abierto en \emph{todos} los puntos de $X_y$, lo que termina de demostrar que \emph{d)} implica \emph{b)}.

El resultado siguiente fue puesto de manifiesto por D.\ Mumford:

\begin{proposition}[14.5.10]
\label{IV.14.5.10}
\oldpage[IV-3]{220}
Sea $Y$ un preesquema noetheriano y $f:X\to Y$ un morfismo universalmente abierto, sobreyectivo y localmente de tipo finito.
Existe entonces un morfismo finito sobreyectivo $g:Y'\to Y$ tal que, si ponemos $X'=X\times_Y Y'$ y $f'=f_{(Y')}:X'\to Y'$, todo punto $y'\in Y'$ admite un entorno abierto $U'$ tal que existe una $U'$-sección de $f'^{-1}(U')$.
\end{proposition}

Demostraremos la proposición en varias etapas.

\emph{I) Reducción al caso en que $Y$ es íntegro.}
Si hemos demostrado la proposición para cada uno de los subpreesquemas reducidos $Y_i$ cuyo espacio subyacente es una componente irreducible de $Y$, y para las imágenes inversas $f^{-1}(Y_i)$, es claro que el preesquema $Y'$ que es la \emph{suma} de los $Y'_i$ correspondientes satisface lo requerido.
Podemos suponer, por tanto, que $Y$ es \emph{íntegro}, y nos limitaremos a este caso en lo que sigue.
Además, en la conclusión podremos tomar también $Y'$ íntegro (tras reemplazarlo, si es necesario, por una componente irreducible adecuada).

\emph{II) Carácter local sobre $Y$.}
Mostraremos que, si $Y$ puede recubrirse por un número finito de abiertos $U_j$ tales que, para todo $j$, la conclusión de la proposición se cumple para el morfismo $f^{-1}(U_j)\to U_j$, restricción de $f$, entonces se cumple también para $f$.
En efecto, podemos suponer evidentemente que los $U_j$ son afines, de modo que $U_j=\Spec(A_j)$, donde $A_j$ es un anillo noetheriano íntegro cuyo cuerpo de fracciones $K=R(Y)$ es el cuerpo de funciones racionales sobre $Y$.
Para todo $j$, existe por hipótesis una $A_j$-álgebra finita e íntegra $A'_j$ tal que el homomorfismo $A_j\to A'_j$ es inyectivo \sref[I]{I.1.2.7} y que el morfismo correspondiente $g_j:U'_j=\Spec(A'_j)\to\Spec(A_j)=U_j$ satisface las condiciones del enunciado (para $U_j$ y $f^{-1}(U_j)$).
Sea $K'$ una extensión finita de $K$ que contiene los cuerpos de fracciones de todas las $A'_j$ (que son extensiones finitas de $K$).
Consideremos la normalización $Y''$ de $Y$ en $K'$ \sref[II]{II.6.3.8}; es de la forma $\Spec(\sh{B})$, donde $\sh{B}$ es una $\sh{O}_Y$-álgebra entera cuasicoherente, cierre entero de $\sh{O}_Y$ en $K'$ \sref[II]{II.6.3.4}.
Estas definiciones muestran que, para todo $j$,
\[
 A'_j=\Gamma(U'_j,\sh{O}_{U'_j})
     =\Gamma\bigl(U_j,(g_j)_*(\sh{O}_{U'_j})\bigr)
\]
se identifica con una sub-$A_j$-álgebra finita de $\Gamma(U_j,\sh{B})$; dicho de otro modo,
\[
 (g_j)_*(\sh{O}_{U'_j})=\sh{A}'_j
\]
es una $\sh{O}_{U_j}$-álgebra coherente, subálgebra de $\sh{B}|U_j$.
Existe, por consiguiente, un sub-$\sh{O}_Y$-módulo coherente $\sh{C}'_j$ de $\sh{B}$ tal que $\sh{C}'_j|U_j=\sh{A}'_j$ \sref[I]{I.9.4.7}.
Si ponemos $\sh{C}'=\sum_j\sh{C}'_j$, la sub-$\sh{O}_Y$-álgebra $\sh{B}'$ de $\sh{B}$ generada por $\sh{C}'$ es coherente, pues $\sh{B}$ es una $\sh{O}_Y$-álgebra entera.
Mostremos ahora que $Y'=\Spec(\sh{B}')$ satisface lo requerido.
En efecto, es claro que el morfismo $g:Y'\to Y$ es finito y sobreyectivo y que $Y'$ es íntegro; además, para todo $j$, $\Gamma(U_j,\sh{B}')$ es una $A'_j$-álgebra finita.
Dicho de otro modo, la restricción $g^{-1}(U_j)\to U_j$ de $g$ se factoriza en la forma
\[
 g^{-1}(U_j)\longrightarrow U'_j\longrightarrow U_j,
\]
y, como la existencia local de secciones es estable por cambio de base, esto demuestra nuestra afirmación, pues todo $y\in Y$ pertenece a algún $U_j$.

\emph{III) Reducción al caso en que $Y$ es íntegro, local y geométricamente unirramificado.}
Supongamos primero que la proposición se ha demostrado cuando $Y$ es íntegro y \emph{local} (con $Y'$ íntegro), y mostremos que es válida cuando $Y$ es afín, íntegro y noetheriano.
En efecto, en virtud de la reducción II), basta demostrar que, para todo punto $y\in Y$, la proposición se satisface para un entorno abierto afín $V$ de $y$ en $Y$.
Sea $Y=\Spec(A)$, $\mathfrak{p}=\mathfrak{j}_y$, sea $Y_1=\Spec(A_\mathfrak{p})$, y pongamos $X_1=X\times_Y Y_1$.
Por hipótesis, existe un morfismo finito sobreyectivo $g_1:Y'_1\to Y_1$, donde $Y'_1=\Spec(B_1)$ y $B_1$ es una $A_\mathfrak{p}$-álgebra finita e íntegra, por tanto un anillo semilocal, tal que $g_1$ satisface las condiciones del enunciado para $Y_1$ y $X_1$.
Si $y'_j$ ($1\leq j\leq r$) son los puntos cerrados de $Y'_1$, existe entonces un recubrimiento de $Y'_1$ por abiertos $U'_j$ tales que $y'_j\in U'_j$ y existe una $U'_j$-sección $h'_j$ de $X\times_Y U'_j$ ($1\leq j\leq r$).
El $A_\mathfrak{p}$-módulo $B_1$ admite un sistema finito de generadores de la forma $z_i/s$ (con $s\in A-\mathfrak{p}$,
\oldpage[IV-3]{221}
$z_i$ enteros sobre $A$), y puede suponerse (multiplicando $s$, si es necesario, por un elemento de $A$) que estos generadores son elementos del cuerpo de fracciones de $B_1$ enteros sobre $A_s$, de modo que, si $V$ es el abierto afín $D(s)=\Spec(A_s)\subset Y$, $Y'_1$ se identifica con $Y_1\times_V V'$, donde $V'$ es el espectro de la $A_s$-álgebra finita generada por los $z_i/s$; $g:V'\to V$ es, por tanto, un morfismo finito sobreyectivo y $g_1=g_{(Y_1)}$.
Además, aplicando el método de \sref{IV.8.1.2},~$a$), puede suponerse que cada $U'_j$ es la imagen inversa por $p:Y'_1\to V'$ de un abierto $W'_j$ de $V'$, que los $W'_j$ recubren $V'$ \sref{IV.8.3.11}, y que cada sección $h'_j$ es de la forma $(v'_j)_{(Y_1)}$, donde $v'_j$ es una $W'_j$-sección de $X\times_Y W'_j$ \sref{IV.8.8.2},~(i).
Estamos, pues, efectivamente reducidos a demostrar la proposición cuando $Y=\Spec(A)$, donde $A$ es un anillo noetheriano local íntegro.
Sabemos entonces que existe una $A$-álgebra finita e íntegra $B$, con el mismo cuerpo de fracciones que $A$, tal que $A\subset B$ y $\Spec(B)$ es geométricamente unirramificado \sref[0]{0.23.2.5} y \sref{IV.6.15.5}; como el morfismo $\Spec(B)\to\Spec(A)$ es sobreyectivo, puede reemplazarse evidentemente $Y$ por $\Spec(B)$ y $X$ por $X\otimes_A B$ para demostrar la proposición.
Razonando como al comienzo de la reducción III), podemos suponer, por tanto, que $A$ es local e íntegro y que $Y=\Spec(A)$ es geométricamente unirramificado.

\emph{IV) Reducción al caso en que $X$ es íntegro y afín, y $f$ es cuasifinito, sobreyectivo, birracional y universalmente abierto.}
Supongamos, pues, $Y=\Spec(A)$ íntegro, local y geométricamente unirramificado.
Existe entonces un subpreesquema irreducible $X_0$ de $X$ tal que la restricción $f_0:X_0\to Y$ de $f$ es un morfismo cuasifinito y \emph{dominante}, y $f_0(X_0)$ contiene al punto cerrado $y$ de $Y$ \sref{IV.14.5.9}; como $Y$ es geométricamente unirramificado, se sigue de \sref{IV.14.4.1} que $f_0$ sigue siendo \emph{universalmente abierto}.
Además, podemos suponer $X_0$ reducido, por tanto íntegro, y vemos que, reemplazando $X$ por $X_0$, podemos suponer $X$ íntegro, $f$ \emph{cuasifinito y dominante}, y que $f(X)$ contiene a $y$; como $f$ es \emph{abierto} y todo abierto de $Y$ que contiene a $y$ es igual a $Y$, $f$ es sobreyectivo.

Sean $\xi$, $\eta$ los puntos genéricos de $X$ e $Y$, respectivamente; $\kappa(\xi)$ es, por tanto, una extensión finita de $K=\kappa(\eta)$.
Por \sref{IV.4.6.8}, existe una extensión \emph{finita} $K'$ de $K$ tal que $(\Spec(\kappa(\xi)\otimes_K K'))_{\mathrm{red}}$ es geométricamente reducido sobre $K'$\footnote{La fuente impresa escribe aquí $K$; el cambio de base y la conclusión siguiente exigen $K'$.} y sus componentes irreducibles son geométricamente irreducibles; se sigue de \sref{IV.4.5.9} y \sref{IV.4.6.1} que los cuerpos residuales de $\Spec(\kappa(\xi)\otimes_K K')$ son extensiones finitas de $K'$ primarias y separables, y son, por tanto, iguales a $K'$.
Aplicando una vez más \sref[0]{0.23.2.5} y \sref{IV.6.15.5}, existe una $A$-álgebra finita e íntegra $B$, que tiene a $K'$ por cuerpo de fracciones, contiene a $A$ y es tal que $\Spec(B)$ es geométricamente unirramificado; como el morfismo $\Spec(B)\to\Spec(A)$ es sobreyectivo, podemos reemplazar de nuevo $Y$ por $\Spec(B)$ y $X$ por $(X\otimes_A B)_{\mathrm{red}}$ para demostrar la proposición; la reducción III) permite entonces suponer una vez más que $A$ es local e íntegro y que $Y=\Spec(A)$ es geométricamente unirramificado, de punto cerrado $y$.
Además, $f$ sigue siendo un morfismo cuasifinito, \emph{sobreyectivo} y \emph{universalmente abierto}; cada una de las componentes irreducibles $X_i$ de $X$ ($1\leq i\leq m$) domina a $Y$ \sref{IV.1.10.4} y es \emph{birracional} sobre $Y$.
Sea, pues, $x$ un punto de $f^{-1}(y)$, y sea $X_i$ una componente irreducible de $X$ que contiene a $x$; denotemos también por $X_i$ el subpreesquema reducido (por tanto, íntegro) de $X$ cuyo espacio subyacente es $X_i$, y sea $f_i:X_i\to Y$ la restricción de $f$.
Aplicando una vez más \sref{IV.14.4.1} al morfismo cuasifinito y dominante $f_i$, vemos que $f_i$ es universalmente abierto; si $U$ es un abierto afín de $X_i$ que contiene a $x$, $f_i(U)$ es, por tanto,
\oldpage[IV-3]{222}
abierto en $Y$ y contiene a $y$, luego es igual a $Y$.
Podemos, pues, reemplazar $X$ por $U$ para demostrar la proposición.

\emph{V) Fin de la demostración.}
Suponemos, pues, $Y$ local, íntegro y geométricamente unirramificado, $X$ íntegro y afín, y $f$ cuasifinito, sobreyectivo, birracional y universalmente abierto; se sigue que $f$ es automáticamente \emph{separado}.
En virtud del teorema principal \sref{IV.8.12.6}, existe una factorización $X\xrightarrow{j}Z\xrightarrow{u}Y$, donde $j$ es una inmersión abierta y $u$ un morfismo \emph{finito}.
Reemplazando además $Z$ por la imagen cerrada de $j$ \sref[I]{I.9.5.10}, podemos suponer que $Z$ es \emph{íntegro}; como $u$ es finito y birracional e $Y$ geométricamente unirramificado, se sigue de \sref[III]{III.4.3.5} y \sref[III]{III.4.3.4} que $u$ (y, por consiguiente, $f$) es un morfismo \emph{radicial}; como $f$ es universalmente abierto y sobreyectivo, es, por tanto, un \emph{homeomorfismo} universal.
En consecuencia, por \sref{IV.2.4.5},~(ii), $f$ es un morfismo \emph{finito}.
Pero entonces se satisfacen las condiciones del enunciado con $Y'$ \emph{íntegro} tomando simplemente $Y'=X$ y $g=f$, siendo la $Y'$-sección el morfismo diagonal $\Delta_f$.

Este resultado puede utilizarse para desarrollar criterios de ``descenso'' para las propiedades de los morfismos universalmente abiertos y sobreyectivos.
Señalemos en particular el criterio siguiente, debido a D.\ Mumford:

\begin{corollary}[14.5.11]
\label{IV.14.5.11}
Sea $Y$ un preesquema localmente noetheriano, $f:X\to Y$ un morfismo localmente de tipo finito y $g:Y'\to Y$ un morfismo localmente de tipo finito, universalmente abierto y sobreyectivo; pongamos $X'=X\times_Y Y'$, $f'=f_{(Y')}:X'\to Y'$.
Para que $f$ sea afín, es necesario y suficiente que $f'$ lo sea.
\end{corollary}

Basta demostrar que la condición es suficiente.
Como la cuestión es local sobre $Y$, podemos suponer que $Y$ es afín, por tanto noetheriano.
En virtud de \sref{IV.14.5.10}, existe un morfismo finito sobreyectivo $h:Y_1\to Y$ tal que, si ponemos $Y'_1=Y'\times_Y Y_1$ y $g'=g_{(Y_1)}:Y'_1\to Y_1$, todo punto de $Y_1$ admite un entorno abierto $U_1$ tal que existe una $U_1$-sección de $g'^{-1}(U_1)$.
Si ponemos $X_1=X\times_Y Y_1$, la proyección canónica $p:X_1\to X$ es un morfismo finito sobreyectivo; por consiguiente, si demostramos que $X_1$ es un esquema afín, se seguirá ante todo que $X$ es cuasicompacto, luego noetheriano, y después que $X$ es afín en virtud del teorema de Chevalley \sref[II]{II.6.7.1}.
Basta, pues, demostrar que el morfismo $f_1=f_{(Y_1)}:X_1\to Y_1$ es afín.
Ahora bien, si ponemos
\[
  X'_1=X_1\times_{Y_1}Y'_1=X'\times_{Y'}Y'_1 \qquad\text{y}\qquad f'_1=(f_1)_{(Y'_1)}=f'_{(Y'_1)}:X'_1\to Y'_1,
\]
$f'_1$ es afín en virtud de la hipótesis.
Estamos, pues, reducidos a demostrar el corolario reemplazando $Y$, $X$, $f$ e $Y'$ por $Y_1$, $X_1$, $f_1$ e $Y'_1$; dicho de otro modo, basta demostrar que $f$ es afín cuando, en el enunciado de \sref{IV.14.5.11}, se supone además que todo punto de $Y$ admite un entorno abierto $U$ tal que existe una $U$-sección de $g^{-1}(U)$.
Como la cuestión es local sobre $Y$, podemos suponer incluso que existe una $Y$-sección $s$ de $Y'$.
Disponemos ahora del lema elemental siguiente (válido en toda categoría que admita productos fibrados):

\begin{lemma}[14.5.11.1]
\label{IV.14.5.11.1}
Sean $f:X\to S$, $g:Y\to S$ dos morfismos, y $p_1:X\times_S Y\to X$, $p_2:X\times_S Y\to Y$ las proyecciones canónicas.
Si $s:S\to Y$ es una sección de $g$, entonces $s'=(1,s\circ f)_S$ es una sección de $p_1$, y $X$, provisto de los morfismos $f:X\to S$ y $s':X\to X\times_S Y$, se identifica con el producto de los $Y$-preesquemas $S$ y $X\times_S Y$ para los morfismos $s:S\to Y$ y $p_2:X\times_S Y\to Y$.
\end{lemma}

\begin{proof}
Este es un caso particular de \sref[I]{I.3.3.11}, donde se reemplaza el diagrama por
\[
\xymatrix{
  X \ar[r]^-{s'} \ar[d]_f & X\times_S Y \ar[r]^-{p_1} \ar[d]^{p_2} & X \ar[d]^f \\
  S \ar[r]_-{s} & Y \ar[r]_-{g} & S
}
\]
\end{proof}

Aplicando este lema con $S$ e $Y$ reemplazados, respectivamente, por $Y$ e $Y'$, se ve que puede escribirse $f=(f')_{(Y)}$ para el cambio de base $s:Y\to Y'$; por tanto, $f$ es afín, puesto que $f'$ lo es.
Q.E.D.

Una variante de este criterio es la siguiente:

\begin{corollary}[14.5.12]
\label{IV.14.5.12}
\oldpage[IV-3]{223}
Con las notaciones e hipótesis generales de \sref{IV.14.5.11}, sea $\sh{L}$ un $\sh{O}_X$-módulo invertible.
Supongamos que existe un subconjunto cerrado $Z$ de $X$, propio sobre $Y$ \sref[II]{II.5.4.10}, y tal que el preesquema inducido por $X$ en el abierto $X-Z$ es normal.
Para que $\sh{L}$ sea amplio con respecto a $f$, es necesario y suficiente que $\sh{L}'=\sh{L}\otimes_{\sh{O}_Y}\sh{O}_{Y'}$ sea amplio con respecto a $f'$.
\end{corollary}

\begin{proof}
Conservando las notaciones de la demostración de \sref{IV.14.5.11}, pongamos $\sh{L}_1=\sh{L}\otimes_{\sh{O}_Y}\sh{O}_{Y_1}$.
Por \sref[III]{III.2.6.2}, bastará demostrar que $\sh{L}_1$ es amplio con respecto a $f\circ p$; pero $f\circ p=h\circ f_1$, y, a la vista de \sref[II]{II.4.6.13},~(v), bastará demostrar que $\sh{L}_1$ es amplio con respecto a $f_1$.
Como la cuestión es local sobre $Y_1$, podemos suponer de nuevo que $g'$ admite una sección $s$.
Si $\sh{L}'_1=\sh{L}'\otimes_{\sh{O}_{Y'}}\sh{O}_{Y'_1}$, puede escribirse entonces, en virtud del lema \sref{IV.14.5.11.1}, $\sh{L}_1=\sh{L}'_1\otimes_{\sh{O}_{Y'_1}}\sh{O}_{Y_1}$ para el cambio de base $s:Y_1\to Y'_1$; la conclusión resulta, por tanto, de dos aplicaciones de \sref[II]{II.4.6.13},~(iii).
\end{proof}
